\documentclass[11pt]{article}
\usepackage[T1]{fontenc}
\usepackage{lmodern,microtype}
\usepackage[margin=1in]{geometry}
\usepackage{amsmath,amssymb,amsthm,mathtools}
\usepackage{booktabs,tabularx,array,enumitem}
\usepackage[dvipsnames]{xcolor}
\usepackage[colorlinks=true,linkcolor=MidnightBlue,citecolor=MidnightBlue,urlcolor=MidnightBlue]{hyperref}
\usepackage{bookmark}
\hypersetup{pdftitle={Localized zero correlations for the Newman theta transform},pdfauthor={Drafted for Edward Baker with OpenAI Codex assistance}}
\setlength{\emergencystretch}{2em}
\setcounter{tocdepth}{1}
\numberwithin{equation}{section}
\newtheorem{theorem}{Theorem}[section]
\newtheorem{proposition}[theorem]{Proposition}
\newtheorem{lemma}[theorem]{Lemma}
\newtheorem{corollary}[theorem]{Corollary}
\theoremstyle{definition}
\newtheorem{definition}[theorem]{Definition}
\theoremstyle{remark}
\newtheorem{remark}[theorem]{Remark}
\newcommand{\R}{\mathbb R}
\newcommand{\C}{\mathbb C}
\newcommand{\dd}{\,\mathrm d}
\newcommand{\ii}{\mathrm i}
\newcommand{\ee}{\mathrm e}
\newcommand{\supp}{\operatorname{supp}}
\newcommand{\diag}{\operatorname{diag}}
\newcommand{\dist}{\operatorname{dist}}
\newcommand{\sgn}{\operatorname{sgn}}
\newcommand{\abs}[1]{\left|#1\right|}
\newcommand{\norm}[1]{\left\|#1\right\|}
\newcolumntype{Y}{>{\raggedright\arraybackslash}X}
\title{Localized zero correlations for the Newman theta transform\\[4pt]
\large Source approximation, collision-time regularity, and finite acquisition}
\author{Drafted for Edward Baker\\\small With substantial OpenAI Codex LLM assistance}
\date{9 November 2026\\\small Working analytical draft; internal assessment levels stated in Appendix~\ref{app:status}}
\begin{document}
\maketitle
\begin{abstract}
We develop localized factorial correlations of the complete Newman theta
zero divisor, retaining real and nonreal locations, multiplicity labels,
and a declared law for the observation center. Gaussian observation
smoothing yields a full hierarchy with explicit normalization and self
corrections. At time zero, logarithmic stability transfers complete-source
accuracy to the entire marked pair function. A polynomial logarithm
regularization then specifies the same-center mixed moments needed to
acquire a canonical profile. For a fixed compact physical window, pairing
the zero forces removes the reciprocal-gap singularity and gives a
Lipschitz modulus through every multiple-zero event, with a complete
exterior error bound. For a Gaussian physical window, exact center and
observation integration produces a separately marked pair projection.
We prove two finite acquisition estimates for it: a second-order centroid
error, valid without internal gap assumptions, and a direct cell-range
bound whose uncertainty is summable by Gaussian separation decay. These
are approximation, regularity, and conditional acquisition results; a
canonical joint profile and the projected scalar's final evaluation
remain outstanding. Finite computations appear only as supporting
applications, with their incomplete gates stated explicitly.
\end{abstract}
\begingroup
\small
\makeatletter
\renewcommand{\l@section}{\@dottedtocline{1}{0em}{1.6em}}
\makeatother
\tableofcontents
\endgroup
\clearpage
\section{The source, the zeros, and the analytical objective}\label{sec:framework}

The primary object is a function of two or more observed zero locations.
It is obtained by marginalizing a full labelled zero configuration, with
spatial localization and Gaussian smoothing made explicit. This retains
ordinary real-zero spacings as well as information on possible nonreal
zeros. A scalar transverse detector is one projection of this object;
it is not the object itself.

The development has three levels. Section~\ref{sec:gaussian} constructs
full factorial densities and proves their exact consistency relations.
Section~\ref{sec:centers} introduces separation coordinates, unfolding,
and a declared law for the observation center. Sections~\ref{sec:logarithms}--\ref{sec:full-pair}
then connect these functions to the complete canonical source, culminating
in the uniform pair approximation of Theorem~\ref{thm:full-pair}.
The identity for the center covariance in Proposition~\ref{prop:covariance}
identifies a further source-specific analytical question. An approximation
bound makes the correlation function accessible to prescribed accuracy;
an estimate of its shape or covariance requires additional analysis.

The new marginal, smoothing and two-sided transfer derivations have
proof-author checking only. They use the local logarithm and one-point
approximation accepted in the project's internal E02 assessment.
Earlier scalar results retain their individual assessment levels and
appear in the appendices. Appendix~\ref{app:status} records these
boundaries. The manuscript asserts neither a determinantal law for
canonical zeros nor a new nonreal-zero exclusion.

\subsection{Canonical normalization and heat equation}

Let
\begin{equation}\label{eq:source}
 r_u=\pi\ee^{4u},\qquad
 \Phi(u)=\ee^u\sum_{\ell\ge1}
       (2\ell^4r_u^2-3\ell^2r_u)\ee^{-\ell^2r_u},
 \qquad
 H_\lambda(z)=\int_0^\infty \ee^{\lambda u^2}\Phi(u)\cos(zu)\dd u .
\end{equation}
The arithmetic coefficients and the complete half-line integral are part
of the definition. For $u\ge0$, $r_u\ge\pi$, so each displayed summand
is positive. More useful for analysis is the bound
\begin{equation}\label{eq:source-basic-majorant}
 |\Phi(u)|\le
 5\ee^u r_u^2\ee^{-r_u}
 \sum_{\ell\ge1}\ell^4\ee^{-\pi(\ell^2-1)}.
\end{equation}
Indeed $3\ell^2r_u\le3\ell^4r_u^2$ and
$\ee^{-(\ell^2-1)r_u}\le\ee^{-\pi(\ell^2-1)}$.
The super-exponential factor $\ee^{-\pi\ee^{4u}}$ dominates every
factor arising from compact sets of $\lambda,z$ and finitely many
derivatives. Differentiation under the integral is therefore justified:
\begin{equation}\label{eq:heat}
 \partial_\lambda H_\lambda=-\partial_z^2H_\lambda .
\end{equation}
For real $\lambda$, $H_\lambda$ is real entire and even. Thus its zeros
are invariant under conjugation and negation, including multiplicities.
Our quantitative approximation theorems below concern $\lambda=0$.
Equation~\eqref{eq:heat} alone does not extend their constants to other
times.

For completeness the normalization can be recovered directly from
the classical theta Mellin identity. Write
$\psi(t)=\sum_{\ell\ge1}\ee^{-\pi\ell^2t}$ and
$\theta(t)=1+2\psi(t)$. Poisson summation for the Gaussian gives
$\theta(t)=t^{-1/2}\theta(1/t)$. Splitting the Mellin integral at one,
and then continuing analytically, gives
\begin{equation}\label{eq:xi-theta}
 \xi(s)=\frac12+\frac{s(s-1)}2
 \int_1^\infty \psi(t)
   \bigl(t^{s/2-1}+t^{(1-s)/2-1}\bigr)\dd t ,
 \qquad
 \xi(s)=\frac12s(s-1)\pi^{-s/2}\Gamma(s/2)\zeta(s).
\end{equation}
To see the constants, initially $\Re s>1$ the full Mellin integral of
$\psi$ is $\pi^{-s/2}\Gamma(s/2)\zeta(s)$. On $(0,1)$ substitute the
theta reflection identity. Its elementary terms integrate to
$1/(s-1)-1/s=1/[s(s-1)]$; the remaining integral is the reflected
half of \eqref{eq:xi-theta}. This proves the formula and its entire
continuation.

Set $s=(1+\ii z)/2$. Under $t=\ee^{4u}$, the cosine in
\eqref{eq:source} splits into the powers corresponding to $s$ and $1-s$,
and
\[
 H_0(z)=\frac18\sum_{\sigma\in\{s,1-s\}}
   \int_1^\infty
   \left(2t^{\sigma/2+1}\psi''(t)
                +3t^{\sigma/2}\psi'(t)\right)\dd t .
\]
Twice integrating by parts gives, for either $\sigma$,
\[
 -2\psi'(1)-(1-\sigma)\psi(1)
 +\frac{\sigma(\sigma-1)}2
       \int_1^\infty t^{\sigma/2-1}\psi(t)\dd t .
\]
Differentiating theta reflection at one yields
$\psi'(1)=-(1+2\psi(1))/8$. The sum of the endpoint terms is
$1/2$, so comparison with \eqref{eq:xi-theta} proves
\begin{equation}\label{eq:xi-normalization}
 H_0(z)=\frac18\xi\!\left(\frac{1+\ii z}{2}\right)
       =\frac18\xi\!\left(\frac{1-\ii z}{2}\right).
\end{equation}
The second equality is $\xi(s)=\xi(1-s)$. A nontrivial zeta zero
$\rho=\beta+\ii\gamma$ is represented by
\begin{equation}\label{eq:coordinates}
 z=x+\ii y=2\gamma+\ii(1-2\beta).
\end{equation}
The classical zero-free half-planes for $\zeta$, obtained from its Euler
product and functional equation, put these zeros in $|y|\le1$; nothing
below requires strictness at the strip boundary. Reflection across the
critical line is conjugation in the $z$ coordinate.

\subsection{Counting measures and factorial marginals}

For a nonzero holomorphic function $V$, its local zero measure is
\begin{equation}\label{eq:zero-measure}
 \nu_V=\sum_{V(a)=0}m_a\delta_a
       =\frac1{2\pi}\Delta\log|V|
\end{equation}
in distributions. To verify the formula locally, factor
$V(z)=(z-a)^{m_a}G(z)$ with $G$ nonvanishing. Then
$\log|G|$ is harmonic, and integration over a punctured disk, followed
by the divergence theorem, gives
$\Delta\log|z-a|=2\pi\delta_a$. A locally finite sum proves
\eqref{eq:zero-measure}.

On any compact observation region, give each multiplicity copy its own
label, say $z_1,\ldots,z_N$. The ordered factorial pair measure is
\begin{equation}\label{eq:factorial}
 \alpha^{(2)}
 =\sum_{j\ne k}\delta_{(z_j,z_k)}
 =\nu\otimes\nu-\diag_*\nu .
\end{equation}
At a zero of multiplicity $m$, the mass at $(z,z)$ is $m(m-1)$.
Subtracting the entire geometric diagonal would incorrectly erase
different labels at the same location.

For $\varphi\in C_c^{1,1}$ the weak formula reads
\begin{equation}\label{eq:weak-onepoint}
 \nu_V(\varphi)=\frac1{2\pi}
       \int\log|V|\,\Delta\varphi\dd A .
\end{equation}
Here and below weak derivatives are used almost everywhere.
Mollification inside a compact support buffer proves this extension
of the smooth formula: the tests converge uniformly, the Laplacians
remain uniformly bounded and converge almost everywhere, and
$\log|V|\in L^1_{\mathrm{loc}}$. These facts also show that every
holomorphic zero-free multiplier contributes zero to the pairing.

\section{From a full configuration to Gaussian-smoothed marginals}\label{sec:gaussian}

We first fix a finite labelled multiset $\mathbf z=(z_1,\ldots,z_N)$.
Labels, rather than geometric positions, distinguish multiplicity copies.
All densities in this section use planar Lebesgue measure $\dd A$ in
each coordinate. The construction is an observation law for a given
configuration, not a hypothesis that canonical zeros are independent.
The factorial-measure convention agrees with the standard correlation
framework \cite[Section 2]{Borodin2009DPP}, extended here to labelled
multiple zeros before smoothing.

\subsection{The full joint density and its marginals}

For $\epsilon>0$ define the normalized planar Gaussian
\begin{equation}\label{eq:gaussian}
 G_\epsilon(u)=\frac1{2\pi\epsilon^2}
     \exp\!\left(-\frac{|u|^2}{2\epsilon^2}\right),
 \qquad \int_\C G_\epsilon\dd A=1.
\end{equation}
Randomly permute the labels and perturb each observed position by an
independent Gaussian with this density. The full joint probability
density is
\begin{equation}\label{eq:full-joint}
 p_{N,\epsilon}(u_1,\ldots,u_N)
 =\frac1{N!}\sum_{\sigma\in\mathfrak S_N}
      \prod_{a=1}^N G_\epsilon(u_a-z_{\sigma(a)}).
\end{equation}
Write $(N)_k=N!/(N-k)!$ for $1\le k\le N$.
Integrating the last $N-k$ variables gives
\begin{equation}\label{eq:joint-marginal}
 p^{(k)}_{N,\epsilon}(\mathbf u)
 =\frac1{(N)_k}\sum_{j_1,\ldots,j_k\ \mathrm{distinct}}
            \prod_{a=1}^k G_\epsilon(u_a-z_{j_a}).
\end{equation}
Indeed the integrated factors each have mass one and each ordered
$k$-tuple occurs $(N-k)!$ times. The associated factorial correlation
function is $R_{k,\epsilon}=(N)_k p^{(k)}_{N,\epsilon}$. Consequently
\begin{equation}\label{eq:fixed-N-consistency}
 \int_{\C^k}R_{k,\epsilon}\dd A^{\otimes k}=(N)_k,
 \qquad
 \int_\C R_{k,\epsilon}(\mathbf u)\dd A(u_k)
       =(N-k+1)R_{k-1,\epsilon}(u_1,\ldots,u_{k-1}).
\end{equation}
These are genuine marginals of a full joint law. They are also
Gaussian convolutions of the deterministic factorial zero measures.
Neither interpretation assigns a new probability distribution to the
unsmoothed canonical zeros.

\subsection{Smooth spatial localization and weighted normalization}

Choose a compactly supported $C^2$ function $0\le W\le1$ on $\C$.
For the zeros of a nonzero entire function only finitely many weights
$w_j=W(z_j)$ are nonzero. Our primary localized objects are
\begin{align}
 \alpha_k^W&=\sum_{j_1,\ldots,j_k\ \mathrm{distinct}}
       \left(\prod_{a=1}^k w_{j_a}\right)
       \delta_{(z_{j_1},\ldots,z_{j_k})},\label{eq:weighted-factorial}\\
 R_{k,\epsilon}^W(\mathbf u)
   &=\int\prod_{a=1}^kG_\epsilon(u_a-z_a)\dd\alpha_k^W(\mathbf z)
     =\sum_{j_1,\ldots,j_k\ \mathrm{distinct}}
          \prod_{a=1}^k w_{j_a}G_\epsilon(u_a-z_{j_a}).\label{eq:full-correlation}
\end{align}
Set $R_{0,\epsilon}^W=1$ and $M_W=\sum_jw_j$.
The window acts on the true zero locations, before the observation
noise. Thus $R_{k,\epsilon}^W$ has Gaussian tails in observed coordinates
even though its source test is compactly supported.

There is a consistent point-process interpretation: retain label $j$
independently with probability $w_j$, and then perturb each retained
location as above. If $M$ is the random retained count, conditioning on
$M=m$ and uniformly labelling the retained observations gives symmetric
probability marginals $p_m^{(k)}$. Their mixture satisfies
\begin{equation}\label{eq:random-N-consistency}
 R_{k,\epsilon}^W
    =\sum_{m\ge k}(m)_k\Pr(M=m)p_m^{(k)},\qquad
 Z_k:=\int R_{k,\epsilon}^W=k!e_k(w_1,\ldots,w_N)\le M_W^k.
\end{equation}
The equality follows by conditioning each ordered retained tuple;
the inequality restores the tuples with repeated labels. If $Z_k>0$,
$R_{k,\epsilon}^W/Z_k$ is a probability density for a weighted ordered
tuple. These separately normalized densities need not be mutually
consistent marginals of a fixed-size law: integrating the last coordinate
in \eqref{eq:full-correlation} leaves the factor
$M_W-\sum_{a<k}w_{j_a}$ for each remaining tuple, not a constant.
The fixed-$N$ rule \eqref{eq:fixed-N-consistency} applies only to the
unweighted retained set, or conditionally to a declared fixed count.

This distinction also separates localization from reweighting an
ensemble joint density by $\prod_jW(z_j)$ and renormalizing. That latter
operation changes the configuration law and its partition function.
Here the canonical configuration is fixed; a Vandermonde ensemble would
be a separate model, not a consequence of the spatial window.

\begin{proposition}[The weighted pair and its self-label correction]\label{prop:weighted-pair}
Define
\begin{equation}\label{eq:r-and-d}
 r_\epsilon^W(u)=\sum_j w_jG_\epsilon(u-z_j),\qquad
 d_\epsilon^W(u,v)=\sum_jw_j^2G_\epsilon(u-z_j)G_\epsilon(v-z_j).
\end{equation}
Then the full two-zero correlation function is
\begin{equation}\label{eq:pair-rd}
 R_{2,\epsilon}^W(u,v)
       =r_\epsilon^W(u)r_\epsilon^W(v)-d_\epsilon^W(u,v),
 \qquad Z_2=M_W^2-\sum_jw_j^2.
\end{equation}
At a location with $m$ multiplicity copies, its coincident-label-group
contribution is $m(m-1)W(z)^2G_\epsilon(u-z)G_\epsilon(v-z)$.
\end{proposition}
\begin{proof}
Expand the product of the two sums defining $r$. Its $j=k$ terms
are exactly $d$, with weight $w_j^2$. The remaining terms are ordered
distinct labels, including distinct labels at the same position.
Integrating both Gaussian factors proves the mass formula.
\end{proof}

The square on $w_j$ in $d$ is essential. The ordinary second moment
of the randomly thinned counting measure instead has a diagonal
expectation weighted by $w_j$, because a retained label is counted once.
Equation~\eqref{eq:pair-rd} subtracts the repeated index in a product
of first intensities, which has weight $w_j^2$.

\subsection{A generating functional for all orders}

For a bounded complex test $f$, let
\begin{equation}\label{eq:full-pgfl}
 \mathcal G_{W,\epsilon}[f]
  =\prod_j\left(1+w_j\int_\C f(u)G_\epsilon(u-z_j)\dd A(u)\right).
\end{equation}
This is the expectation of $\prod_{u\ \mathrm{observed}}(1+f(u))$ under
the thinning-and-smoothing law. Expanding the finite product yields
\begin{equation}\label{eq:pgfl-expansion}
 \mathcal G_{W,\epsilon}[f]
  =\sum_{k=0}^N\frac1{k!}
      \int_{\C^k}\prod_{a=1}^k f(u_a)
                       R_{k,\epsilon}^W(\mathbf u)\dd A^{\otimes k}.
\end{equation}
If $\|f\|_\infty\le q$, the absolute tail after order $K$ is at most
\begin{equation}\label{eq:order-tail-full}
 \sum_{k>K}\frac{(qM_W)^k}{k!}
       \le \ee^{qM_W}\frac{(qM_W)^{K+1}}{(K+1)!}.
\end{equation}
For the last inequality put $k=K+1+j$ and use
$(K+1+j)!\ge(K+1)!j!$. This bound pays the order truncation
without identifying different orders or postulating a closed hierarchy.

\subsection{Smoothing flow, bias, and localization tails}

Put $\tau=\epsilon^2/2$. Direct differentiation of
$G_{\sqrt{2\tau}}=(4\pi\tau)^{-1}\ee^{-|u|^2/(4\tau)}$ gives
$\partial_\tau G_{\sqrt{2\tau}}=\Delta G_{\sqrt{2\tau}}$.
The finite sums therefore satisfy
\begin{equation}\label{eq:smoothing-flow}
 \partial_\tau R_{k,\sqrt{2\tau}}^W
    =\sum_{a=1}^k\Delta_{u_a}R_{k,\sqrt{2\tau}}^W,
 \qquad R_{k,\sqrt{2\tau}}^W\dd A^{\otimes k}
      \rightharpoonup\alpha_k^W\quad(\tau\downarrow0).
\end{equation}
Weak convergence follows by testing each Gaussian approximate identity.
This is an auxiliary observation-smoothing time. It is distinct from
Newman time $\lambda$ in \eqref{eq:heat}.

For $F\in C_b^2(\C^k)$, with $\mathcal L_k=\sum_a\Delta_{u_a}$,
the Gaussian semigroup and its sup-norm contraction give
\begin{equation}\label{eq:smoothing-bias}
 \left|\int F R_{k,\epsilon}^W\dd A^{\otimes k}
                       -\alpha_k^W(F)\right|
 \le\frac{\epsilon^2}{2}\,Z_k\|\mathcal L_kF\|_\infty.
\end{equation}
Indeed $P_\tau F-F=\int_0^\tau P_t\mathcal L_kF\dd t$; integrate
against the positive measure $\alpha_k^W$. This is a weak bias estimate,
not uniform convergence of densities to an atomic measure.

For a countable weighted configuration with $M_W<\infty$, all the
preceding nonnegative sums converge. If a truncation omits total weight
$T$, let $R_{k,\epsilon}^{W,\mathrm{tr}}$ denote the retained sum and put
$A_\epsilon=(2\pi\epsilon^2)^{-1}$. Then
\begin{equation}\label{eq:spatial-tail}
 \begin{split}
 0&\le R_{k,\epsilon}^W-R_{k,\epsilon}^{W,\mathrm{tr}}
          \le k A_\epsilon^kM_W^{k-1}T,\\
 \|R_{k,\epsilon}^W-R_{k,\epsilon}^{W,\mathrm{tr}}\|_{L^1}
          &\le kM_W^{k-1}T.
 \end{split}
\end{equation}
Every omitted tuple has an omitted label in at least one of $k$ slots;
sum over the choice of slot and restore repeated labels to bound the
remaining $k-1$ sums by $M_W^{k-1}$. This also proves the $L^1$ claim
after integrating the Gaussians. Section~\ref{sec:gaussian-window}
pays the corresponding canonical mass and tail estimates for an
untruncated Gaussian physical window. The source theorem in
Section~\ref{sec:full-pair} uses a compact physical window and Gaussian
observation smoothing.

\section{Separation, unfolding, and the center-averaged law}\label{sec:centers}

The pair function in \eqref{eq:full-correlation} has four real arguments.
Each reduction of that information should be stated as an explicit
coordinate change or marginalization. This is especially relevant when
comparing a complex-zero statistic with a real-line spacing law.

\subsection{Center and separation coordinates}

For $u,v\in\C$ set
\begin{equation}\label{eq:center-separation}
 c=\frac{\Re u+\Re v}{2},\quad s=\Re u-\Re v,
 \quad a=\Im u,\quad b=\Im v,
 \qquad
 \mathcal C_{2,\epsilon}^W(c,s,a,b)
    =R_{2,\epsilon}^W(c+s/2+\ii a,c-s/2+\ii b).
\end{equation}
The absolute Jacobian is one. Thus integrating $c$ gives a separation
density with two transverse marks; subsequently integrating $a,b$
gives a horizontal separation density. Neither is the full function.
Dividing any of these by $Z_2>0$ gives its ordered-pair probability
density. The ratio $R_2(u,v)/(r(u)r(v))$, where the denominator is
positive, is a different normalization and does not generally have
total mass one.

A second useful projection integrates the full complex center:
\begin{proposition}[Gaussian separation marginal]\label{prop:separation}
For $s\in\C$,
\begin{equation}\label{eq:separation-gaussian}
 \mathcal D_{2,\epsilon}^W(s)
 :=\int_\C R_{2,\epsilon}^W(c+s/2,c-s/2)\dd A(c)
 =\sum_{j\ne k}w_jw_k
                  G_{\sqrt2\epsilon}(s-(z_j-z_k)).
\end{equation}
Its total mass is $Z_2$.
\end{proposition}
\begin{proof}
For a fixed pair, substitute the two Gaussian factors in the integral.
Completing the square in $c$ gives the density of the difference of
two independent Gaussians, whose covariance is $2\epsilon^2 I$.
Its normalization is $(4\pi\epsilon^2)^{-1}$. Summing and then
integrating in $s$ proves both claims. The map
$(c,s)\mapsto(c+s/2,c-s/2)$ has absolute Jacobian one in $\R^4$.
\end{proof}
Unlike \eqref{eq:center-separation}, this projection integrates out the
common transverse coordinate. It cannot recover two independent marks.

If all zeros under observation are real, write
$g_\epsilon(t)=(2\pi\epsilon^2)^{-1/2}\ee^{-t^2/(2\epsilon^2)}$.
The exact factorization is
\begin{equation}\label{eq:real-zero-correlation}
 R_{2,\epsilon}^W(u,v)
 =g_\epsilon(\Im u)g_\epsilon(\Im v)
   \sum_{j\ne k}w_jw_k
        g_\epsilon(\Re u-x_j)g_\epsilon(\Re v-x_k).
\end{equation}
The transverse spread here is known observation noise; the horizontal
factor retains the full real-zero spacing information. A transversely
vanishing scalar detector would discard this nonzero pair function.

\subsection{A declared center law and an unfolding map}

Let $X$ range over a compact interval $I\subset[10,\infty)$ with a
specified probability measure $\pi$. Examples are normalized Lebesgue
measure on $I$ and a Gaussian density truncated to $I$ and renormalized.
For each $X$, declare the window $W_X$, smoothing width $\epsilon_X>0$,
and positive horizontal and transverse scales $\ell_X,b_X$. Define
\begin{equation}\label{eq:unfolding}
 T_X(\zeta)=X+\ell_X\Re\zeta+\ii b_X\Im\zeta,
 \quad J_X=\ell_X b_X,\qquad
 \widetilde R_{k,X}(\boldsymbol\zeta)
   =J_X^kR_{k,\epsilon_X}^{W_X}(T_X\zeta_1,\ldots,T_X\zeta_k).
\end{equation}
The Jacobian preserves total factorial mass. It is legitimate to choose
$\ell_X$ as a nominal mean-spacing scale, but that choice does not
prove that the actual local mean density is constant. We smooth first
in physical coordinates, then unfold; the two resulting Gaussian widths
are $\epsilon_X/\ell_X$ and $\epsilon_X/b_X$. Fixed isotropic smoothing
after unfolding is a different convention, implemented by an anisotropic
physical Gaussian.

Average the conditional observation laws over $X$. Their correlation
functions, including all orders, are
\begin{equation}\label{eq:center-average}
 \overline R_k(\boldsymbol\zeta)
       =\int_I\widetilde R_{k,X}(\boldsymbol\zeta)\dd\pi(X).
\end{equation}
This is a fully specified statistical ensemble derived from the
canonical configuration and an external center law. For simplicity,
all parameter choices below are measurable and bounded on $I$, with
positive lower bounds on the scales and smoothing width. A Gaussian
law on the whole real line would require an additional center-tail
bound; truncation is part of the definition here.

\begin{proposition}[The connected center-averaged pair]\label{prop:covariance}
Let $\widetilde r_X(\zeta)=J_Xr_{\epsilon_X}^{W_X}(T_X\zeta)$ and
$\widetilde d_X(\zeta,\eta)=J_X^2d_{\epsilon_X}^{W_X}(T_X\zeta,T_X\eta)$.
Then, with $\overline r=\int_I\widetilde r_X\dd\pi$,
\begin{equation}\label{eq:center-covariance}
 \overline R_2(\zeta,\eta)-\overline r(\zeta)\overline r(\eta)
 =\operatorname{Cov}_{\pi}\bigl(\widetilde r_X(\zeta),
                              \widetilde r_X(\eta)\bigr)
       -\int_I\widetilde d_X(\zeta,\eta)\dd\pi(X).
\end{equation}
Here covariance is the bilinear real covariance of the displayed
real-valued densities.
\end{proposition}
\begin{proof}
Multiply \eqref{eq:pair-rd} by $J_X^2$, substitute $T_X\zeta,T_X\eta$,
and average. Subtract the product of the two averaged first intensities.
The remaining product term is exactly the covariance in
\eqref{eq:center-covariance}.
\end{proof}

At a fixed center the connected function for the auxiliary independent
observations is just $-d$. The full pair function still contains the
configuration's spatial geometry. Averaging over centers adds the
covariance of that geometry. Thus a canonical estimate of the
center covariance is a concrete further analytical target; it is not
provided by independence of the observation noise. A source expression
for this covariance is obtained in Section~\ref{sec:full-pair}.

The CUE calculation in Appendix~\ref{app:ensemble} provides an explicit
comparison law derived from a supplied ensemble. No such ensemble law
is assumed for the canonical divisor. Its center covariance requires
the additional information isolated in Section~\ref{sec:joint-source}.

\section{From holomorphic accuracy to logarithmic accuracy}\label{sec:logarithms}

Uniform approximation of a holomorphic function does not give a
pointwise logarithm comparison near its zeros. The appropriate local
statement uses an integral norm, a nonvanishing anchor, and a larger
disk. Multiplicity is then paid through a zero-count exponent.
The following constants are deliberately explicit rather than
optimized.

Fix once and for all
\begin{equation}\label{eq:radii}
 R=\frac92,\qquad \rho_*=\frac{13}3,\qquad r=\frac{25}6,
 \qquad B=R+r=\frac{26}3,\qquad
 \frac{\rho_*+r}{\rho_*-r}=51 .
\end{equation}
Write $D_t=D(0,t)$ in this section.

\begin{theorem}[Local logarithm comparison]\label{thm:log}
Suppose $f,g$ are holomorphic on a neighborhood of $\overline D_R$ and
\[
 |f|,|g|\le1,\qquad |f(0)|,|g(0)|\ge m,\qquad
 0<m\le1,\qquad \norm{f-g}_{L^\infty(D_R)}\le\varepsilon\le1 .
\]
Set
\begin{equation}\label{eq:local-constants}
 \begin{gathered}
 A=\log(1/m),\qquad
 n=\max\!\left(1,\left\lceil\frac{A}{\log(27/26)}\right\rceil\right),
 \quad \alpha=\frac2{n+2},\\
 C(A,n)=\pi r^2+\pi B^2n^2\ee^{102A/n}2^{2/n}.
 \end{gathered}
\end{equation}
Then, with locally integrable logarithms at zeros,
\begin{equation}\label{eq:log-transfer}
 \int_{D_r}\left|\log|f|-\log|g|\right|\dd A
 \le C(A,n)\varepsilon^\alpha .
\end{equation}
Each function has at most $n$ zeros in $D_{\rho_*}$, counted with
multiplicity. The same statements hold after translating the disks.
\end{theorem}

\begin{proof}
We give the zero count, sublevel estimate and comparison separately.
The classical inputs are Jensen's formula, the maximum principle and
the disk Harnack inequality; see
\cite{Tao2020LocalZeros,Tao2016Harnack}.

\emph{Zero count and factor removal.}
If $a_1,\ldots,a_k$ are the zeros of $f$ in $D_{\rho_*}$, repeated
according to multiplicity, Jensen's formula and $|f|\le1$ give
\[
 k\log(R/\rho_*)\le
 \sum_{|a|<R}\log(R/|a|)
 \le-\log|f(0)|\le A .
\]
The formula follows at a boundary zero by an interior-radius limit.
Thus $k\le n$, and no $a_\ell$ is zero. Remove these factors by
\[
 b_\ell(z)=\frac{R(z-a_\ell)}{R^2-\overline{a_\ell}z},\qquad
 h(z)=\frac{f(z)}{\prod_{\ell=1}^k b_\ell(z)}.
\]
The apparent poles are removable. Each $b_\ell$ has modulus one on
$|z|=R$, so $|h|\le1$ by the maximum principle.
Moreover $h$ is zero-free in $D_{\rho_*}$ and
$|h(0)|\ge|f(0)|\ge m$. Harnack applied to the nonnegative harmonic
function $-\log|h|$ yields
$-\log|h(z)|\le51A$ on $D_r$. Zeros on the circle of radius
$\rho_*$ cause no problem: apply Harnack on smaller circles and
let their radii increase to $\rho_*$.
For $|z|\le r$,
\begin{equation}\label{eq:blaschke-lower}
 |b_\ell(z)|\ge\frac{|z-a_\ell|}{B},\qquad
 |f(z)|\ge\ee^{-51A}\prod_{\ell=1}^k\frac{|z-a_\ell|}{B}.
\end{equation}
The first inequality follows from
$|R^2-\overline{a_\ell}z|\le R(R+r)$; each distance ratio is
at most one.

\emph{Sublevel areas.}
If $q=\ee^{51A}t<1$ and $k>0$, a point with $|f(z)|<t$ must lie
in one of the $k$ disks of radius $Bq^{1/k}$ around the $a_\ell$.
Their union has area at most
$k\pi B^2q^{2/k}\le n\pi B^2q^{2/n}$.
For $k=0$ the sublevel set is empty. If $q\ge1$, the same upper
bound exceeds the area of $D_r$. Hence for every $t>0$,
\begin{equation}\label{eq:sublevel}
 \bigl|\{z\in D_r:|f(z)|<t\}\bigr|
 \le L(A,n)t^{2/n},\qquad
 L(A,n)=n\pi B^2\ee^{102A/n}.
\end{equation}
Layer cake, applied to $\log(t/|f|)_+$, now gives
\begin{equation}\label{eq:layer-cake}
 \int_{\{|f|<t\}\cap D_r}\log\frac{t}{|f|}\dd A
 \le\int_0^\infty L(A,n)(t\ee^{-v})^{2/n}\dd v
 =\frac{\pi B^2n^2}{2}\ee^{102A/n}t^{2/n}.
\end{equation}
The same estimates hold for $g$.

\emph{Comparison.}
For $\varepsilon>0$ choose
$\tau=\varepsilon^{n/(n+2)}$, so $\varepsilon\le\tau$.
Where both moduli are at least $\tau$, the real logarithm is
$1/\tau$-Lipschitz, and the integral of the log difference is at
most $\pi r^2\varepsilon/\tau$.
On the complement the smaller modulus is less than $\tau$ and the
larger is less than $2\tau$. There
\[
 |\log|f|-\log|g||
 \le\log(2\tau/|f|)+\log(2\tau/|g|).
\]
Extend each nonnegative term to its full sublevel set and apply
\eqref{eq:layer-cake} twice. The resulting bound is
\[
 \pi r^2\frac{\varepsilon}{\tau}
 +\pi B^2n^2\ee^{102A/n}(2\tau)^{2/n}
 =C(A,n)\varepsilon^{2/(n+2)}.
\]
When $\varepsilon=0$ the assertion is immediate.
\end{proof}

\begin{corollary}[Compact zero-statistic transfer]\label{cor:weak-transfer}
Under the hypotheses of Theorem~\ref{thm:log}, for any
$\varphi\in C_c^{1,1}(D_r)$,
\begin{equation}\label{eq:weak-transfer}
 |\nu_f(\varphi)-\nu_g(\varphi)|
 \le \frac{\norm{\Delta\varphi}_\infty}{2\pi}
              C(A,n)\varepsilon^{2/(n+2)} .
\end{equation}
No simplicity, matching of roots, or lower bound for the modulus
throughout the observation region is required.
\end{corollary}
\begin{proof}
Subtract \eqref{eq:weak-onepoint} for $f$ and $g$ and apply
\eqref{eq:log-transfer}.
\end{proof}

We will also use an elementary bound on the individual logarithms:
\begin{equation}\label{eq:log-mass}
 \norm{\log|f|}_{L^1(D_r)},\
 \norm{\log|g|}_{L^1(D_r)}
 \le \pi r^2A .
\end{equation}
Both logarithms are nonpositive, while Jensen's formula says that each
circular mean is at least its value at the center, hence at least
$-A$. Integrating those means in the radius proves
\eqref{eq:log-mass}. This estimate concerns the size of the logarithms;
\eqref{eq:log-transfer} concerns their difference and is the useful
small quantity.

\section{Complete-source approximation at time zero}\label{sec:canonical}

We now supply all hypotheses of Theorem~\ref{thm:log} for the canonical
field at a declared height. The off-axis center supplies a safe anchor;
the observation tests still reach the real axis.

\subsection{A zero-free normalization and a safe anchor}

For $X\ge10$ set
\begin{equation}\label{eq:canonical-F}
 \begin{gathered}
 c=X+4\ii,\qquad s(z)=\frac12-\frac{\ii z}2,\\
 F(z)=\frac{s(z)-1}{s(z)}\zeta(s(z))
     =D(z)H_0(z),\qquad
 D(z)=\frac{16\pi^{s(z)/2}}{s(z)^2\Gamma(s(z)/2)} .
 \end{gathered}
\end{equation}
Throughout $\overline D(c,R)$,
\begin{equation}\label{eq:s-range}
 \frac14\le\Re s\le\frac{19}4,\qquad
 |\Re z|\le X+\frac92 .
\end{equation}
The multiplier $D$ is holomorphic and zero-free on a neighborhood
of this disk. The zeta pole at $s=1$ is cancelled in $F$, while
$s=0$ and the zeros of $1/\Gamma(s/2)$ are outside.

For $\sigma=\Re s>0$, partial summation, initially at $\sigma>1$,
and then analytic continuation give
\begin{equation}\label{eq:fractional-part}
 \frac{s-1}{s}\zeta(s)
 =1-(s-1)\int_1^\infty\{v\}v^{-s-1}\dd v .
\end{equation}
Indeed $\zeta(s)=s\int_1^\infty\lfloor v\rfloor v^{-s-1}\dd v$,
and $\lfloor v\rfloor=v-\{v\}$. The last integral in
\eqref{eq:fractional-part} converges absolutely for $\sigma>0$.
On the disk, $|s-1|\le X/2+6$, so
\begin{equation}\label{eq:canonical-upper}
 |F|\le1+\frac{|s-1|}{\sigma}\le2X+25=M_X-1,
 \qquad M_X=2X+26 .
\end{equation}
At $c$, $\Re s=5/2$. The absolutely convergent Euler product
\cite{DLMFEulerProductOC} and $|(s-1)/s|\ge3/5$ imply
\begin{equation}\label{eq:c0}
 |F(c)|\ge c_0,\qquad
 c_0=\frac35\,\frac{\zeta(5)}{\zeta(5/2)}>0 .
\end{equation}
In fact
$|\zeta(s)|\ge\prod_p(1+p^{-5/2})^{-1}
=\zeta(5)/\zeta(5/2)$, and the rational factor follows by comparing
$(3/2)^2+(\Im s)^2$ with $(5/2)^2+(\Im s)^2$.
This anchor uses an Euler product only in its absolute convergence
half-plane.

If $\widehat F$ is holomorphic near the disk and
\begin{equation}\label{eq:delta-requirement}
 \norm{F-\widehat F}_{L^\infty(D(c,R))}\le\delta
 \le\min(1,c_0/2),
\end{equation}
then $f=F/M_X$ and $g=\widehat F/M_X$ satisfy the preceding theorem
with the common constants
\begin{equation}\label{eq:canonical-constants}
 \begin{gathered}
 m_X=\frac{c_0}{2M_X},\quad A_X=\log\frac{2M_X}{c_0},\quad
 n_X=\max\!\left(1,\left\lceil\frac{A_X}{\log(27/26)}\right\rceil\right),\\
 \alpha_X=\frac2{n_X+2},\qquad C_X=C(A_X,n_X).
 \end{gathered}
\end{equation}
Consequently
\begin{equation}\label{eq:EX}
 \int_{D(c,r)}|\log|F|-\log|\widehat F||\dd A
 \le E_X,\qquad E_X=C_X(\delta/M_X)^{\alpha_X}.
\end{equation}
The factors $M_X$ cancel in the logarithm difference.
Here $n_X=O(\log X)$, $C_X=O((\log X)^2)$ and
$\alpha_X\asymp1/\log X$. For example
$A_X/n_X\le\log(27/26)$ makes the exponential in $C_X$
uniformly bounded. No height-independent positive exponent is claimed.

\subsection{Both source tails and the enlarged-disk multiplier}

For integers $J\ge1$ and $U>0$, let $\Phi_J$ be the sum in
\eqref{eq:source} restricted to $\ell\le J$, and define the exact
finite-source field
\begin{equation}\label{eq:finite-source}
 \widehat H_{J,U}(z)=\int_0^U\Phi_J(u)\cos(zu)\dd u,\qquad
 \widehat F=D\widehat H_{J,U}.
\end{equation}
These are exact finite sums and compact integrals; no quadrature or
rounding has been substituted. The real cosine source ensures
conjugation symmetry of $\widehat H_{J,U}$.

Put $b_* =17/2$ and
\begin{align}
 C_2(J)&=\sum_{\ell>J}\ell^4\ee^{-\pi(\ell^2-1)},
 & B_0&=5C_2(0),\label{eq:source-tail-defs}\\
 I_{b_*}&=\int_0^\infty
       \ee^{(b_*+1)u}r_u^2\ee^{-r_u}\dd u,
 &I_{b_*}^{>U}&=\int_U^\infty
       \ee^{(b_*+1)u}r_u^2\ee^{-r_u}\dd u .\notag
\end{align}
On the enlarged disk, $|\Im z|\le b_*$ and
$|\cos(zu)|\le\ee^{b_*u}$.
Splitting the omitted arithmetic terms on $[0,U]$ from the
\emph{complete} radial tail on $(U,\infty)$, and applying
\eqref{eq:source-basic-majorant}, proves
\begin{equation}\label{eq:delta-H}
 \sup_{|\Im z|\le b_*}|H_0(z)-\widehat H_{J,U}(z)|
 \le\delta_H(J,U):=5C_2(J)I_{b_*}+B_0I_{b_*}^{>U}.
\end{equation}
The integral is a half-line cosine integral, so there is no extra
factor of two.

Both tails have explicit majorants. For $\ell\ge J+1$, consecutive
terms in $C_2(J)$ have ratio at most
$16\ee^{-\pi(2J+3)}<1$. Therefore
\begin{equation}\label{eq:arithmetic-tail}
 C_2(J)\le
 \frac{(J+1)^4\ee^{-\pi((J+1)^2-1)}}
      {1-16\ee^{-\pi(2J+3)}}\qquad(J\ge0).
\end{equation}
The substitution $v=\pi\ee^{4u}$, with $\dd u=\dd v/(4v)$,
gives
\begin{equation}\label{eq:radial-tail}
 I_{b_*}^{>U}=\frac{\pi^{-19/8}}4
       \Gamma(35/8,\pi\ee^{4U}),\qquad
 I_{b_*}=\frac{\pi^{-19/8}}4\Gamma(35/8,\pi),
\end{equation}
where $\Gamma(a,t)=\int_t^\infty v^{a-1}\ee^{-v}\dd v$.

An explicit bound for $D$ on the whole enlarged disk is
\begin{equation}\label{eq:multiplier-bound}
 |D(z)|\le\mathcal B_X,\qquad
 \mathcal B_X=
 32\pi^{19/8}
 \exp\!\left(\frac{19}8+
       \frac{\pi^2(361+4V_X^2)}{768}\right),
 \qquad V_X=X+\frac92 .
\end{equation}
To verify it, the reciprocal-gamma product
\cite{DLMFReciprocalGammaOC} gives
\[
 \frac1{\Gamma(w)}
   =w\ee^{\gamma_{\!E}w}\prod_{k\ge1}(1+w/k)\ee^{-w/k},
 \qquad
 |\Gamma(w)^{-1}|
 \le |w|\exp\!\left(\gamma_{\!E}\Re w+
                      \frac{\pi^2|w|^2}{12}\right).
\]
Here $\log|1+q|\le\Re q+|q|^2/2$ follows from
$\log(1+t)\le t$ applied to $|1+q|^2$, and
$\sum k^{-2}=\pi^2/6$. Apply the estimate to $w=s/2$,
use $\gamma_{\!E}\le1$, $|s|\ge1/4$ and
$|s|^2\le361/16+V_X^2/4$. The prefactor in $D$ is then
$8/|s|\le32$, proving \eqref{eq:multiplier-bound}.
This bound is expensive in $X$; it is used for completeness, not
as a practical-cost assertion.

The input error for \eqref{eq:delta-requirement} is consequently
\begin{equation}\label{eq:complete-delta}
 \delta=\mathcal B_X
       \left[5C_2(J)I_{b_*}+B_0I_{b_*}^{>U}\right].
\end{equation}
Both independent tails tend to zero. Thus any positive field-error
target is attained by some finite $J,U$ at every declared $X$.
No endpoint-reflection identity for the truncated source is assumed.


\section{Canonical approximation of the full pair function}\label{sec:full-pair}

The preceding source estimates were obtained on an off-axis disk with
a safe anchor. A full pair function must also retain real zeros and
zeros on both sides of the axis. Folding a transversely vanishing test
does not accomplish this. We instead use two reflected disks and a
smooth partition that counts every real label with total weight one.
This section contains the quantitative result \cite{OCGaussian},
accepted in the internal E07 assessment after the compact-center
qualification below; see Appendix~\ref{app:status}.

\subsection{A two-sided window with an explicit support buffer}

Define the $C^2$ cutoff
\begin{equation}\label{eq:full-cutoff}
 P(t)=\begin{cases}
 1,&t\le0,\\
 1-10t^3+15t^4-6t^5,&0<t<1,\\
 0,&t\ge1.
 \end{cases}
\end{equation}
Its first derivative is $-30t^2(1-t)^2$ on $(0,1)$, so
$|P'|\le15/8$; differentiation once more gives $|P''|\le6$.
For $X\ge10$ and $0<h\le1$ put
\begin{equation}\label{eq:full-window}
 \begin{gathered}
 w_{X,h}(x)=P(2|x-X|/h-1),\qquad \kappa_0(y)=P(|y|-1),\\
 W_X(x+\ii y)=w_{X,h}(x)\kappa_0(y),\qquad
 \theta_+(y)=P(1/2-16y),\quad\theta_-(y)=1-\theta_+(y).
 \end{gathered}
\end{equation}
The absolute values cause no singular derivatives because each cutoff
is constant near their nonsmooth points. The window equals one for
$|x-X|\le h/2$, $|y|\le1$, and vanishes outside
$|x-X|\le h$, $|y|\le2$. For true time-zero zeros the transverse
cutoff is one; for approximating zeros the collar $1<|y|<2$ keeps
its stated weights.

Let $c_\pm=X\pm4\ii$ and retain $R=9/2$, $\rho_*=13/3$, $r=25/6$
from \eqref{eq:radii}. The support of $\theta_+W_X$ has
$y\ge-1/32$ and hence
\begin{equation}\label{eq:two-disk-buffer}
 |z-c_+|^2\le1+(4+1/32)^2
        =\frac{17665}{1024}<\frac{625}{36}=r^2.
\end{equation}
The reflected statement holds for $\theta_-W_X$. Both supports have a
strict buffer inside their disks. At $y=0$, $\theta_+=\theta_-=1/2$.
Thus the partition retains real zeros exactly, without an axis indicator
or a boundary distribution.

Use $F_+=F=D H_0$ and $\widehat F_+=D\widehat H_{J,U}$ from
Section~\ref{sec:canonical}. Define
$F_-(z)=\overline{F_+(\bar z)}$ and
$\widehat F_-(z)=\overline{\widehat F_+(\bar z)}$.
Conjugation symmetry of both entire fields makes these products of the
same fields with reflected, locally zero-free multipliers. Reflection
transports the source error, anchor and log estimate to the lower disk:
\begin{equation}\label{eq:both-log-errors}
 \int_{D(c_\sigma,r)}
       |\log|F_\sigma|-\log|\widehat F_\sigma||\dd A\le E_X,
 \qquad \sigma\in\{+,-\}.
\end{equation}
Each field has at most $n_X$ zero labels in each $\rho_*$ disk, by
the same Jensen bound. Thus at most $M_X^{\rm pop}:=2n_X$ labels
of either field can have nonzero $W_X$ weight. This population bound
is distinct from the normalization constant $M_X=2X+26$.

\begin{lemma}[Two-sided weak transfer]\label{lem:two-sided}
For any compact $C^2$ test $\psi$ supported in the window rectangle,
\begin{equation}\label{eq:two-sided-transfer}
 \left|(\nu_{H_0}-\nu_{\widehat H_{J,U}})(\psi)\right|
 \le\frac{E_X}{2\pi}
       \sum_{\sigma=\pm}\|\Delta(\theta_\sigma\psi)\|_\infty.
\end{equation}
\end{lemma}
\begin{proof}
Write $\psi=\theta_+\psi+\theta_-\psi$. Each summand is supported
strictly within its inner disk. Apply \eqref{eq:weak-onepoint} there
to the difference of zero measures, remove each zero-free multiplier,
and use \eqref{eq:both-log-errors}. This argument treats all labels,
including coincident labels and labels on the axis, without simplicity
or zero separation assumptions.
\end{proof}

\subsection{Uniform Gaussian derivative costs}

All suprema in the following estimates are over the whole relevant
Euclidean space. Direct radial differentiation of \eqref{eq:gaussian}
gives
\begin{equation}\label{eq:gaussian-costs}
 \|G_\epsilon\|_\infty=A_\epsilon=\frac1{2\pi\epsilon^2},\quad
 \|\nabla G_\epsilon\|_\infty=B_\epsilon
       =\frac{A_\epsilon}{\sqrt{\ee}\epsilon},\quad
 \|\Delta G_\epsilon\|_\infty=C_\epsilon
       =\frac{2A_\epsilon}{\epsilon^2}.
\end{equation}
For the gradient, maximize $t\ee^{-t^2/2}$ at $t=1$; for the
Laplacian use $|t^2-2|\ee^{-t^2/2}\le2$ for $t\ge0$.

The derivative bounds for the window and partition are
$|w'|\le15/(4h)$, $|w''|\le24/h^2$,
$|\theta_\sigma'|\le30$, $|\theta_\sigma''|\le1536$.
The derivatives of $\theta_\sigma$ are supported in $|y|<1/32$,
where $\kappa_0=1$; the derivatives of $\kappa_0$ are supported in
$1<|y|<2$, where $\theta_\sigma$ is constant. This disjointness
removes their cross terms. With
$a_\sigma=\theta_\sigma W_X$ and
$b_\sigma=\theta_\sigma W_X^2$, we obtain
\begin{equation}\label{eq:window-costs}
 \begin{aligned}
 |a_\sigma|,|b_\sigma|&\le1,\\
 \|\nabla a_\sigma\|_\infty&\le L_h=15/(4h)+30,&
 \|\Delta a_\sigma\|_\infty&\le D_h=24/h^2+1536,\\
 \|\nabla b_\sigma\|_\infty&\le L_h^{(2)}=15/(2h)+30,&
 \|\Delta b_\sigma\|_\infty&\le D_h^{(2)}=609/(8h^2)+1536.
 \end{aligned}
\end{equation}
For example $|(w^2)''|\le2|w''|+2|w'|^2\le609/(8h^2)$.
Likewise $|(\kappa_0^2)''|\le609/32<1536$, so the vertical
bound is dominated by the partition transition. Bounding the Euclidean
gradient by the sum of its coordinate bounds proves the stated $L$'s.

For brevity write $A=A_\epsilon$, $B=B_\epsilon$, $C=C_\epsilon$,
and define
\begin{equation}\label{eq:UV-costs}
 \begin{split}
 U_{h,\epsilon}&=D_hA+2L_hB+C,\\
 V_{h,\epsilon}&=D_h^{(2)}A^2+4L_h^{(2)}AB+2AC+2B^2.
 \end{split}
\end{equation}
The product rule gives, uniformly in $u,v\in\C$,
\begin{align}
 \|\Delta_z[a_\sigma(z)G_\epsilon(u-z)]\|_\infty
       &\le U_{h,\epsilon},\label{eq:single-test-cost}\\
 \|\Delta_z[b_\sigma(z)G_\epsilon(u-z)G_\epsilon(v-z)]\|_\infty
       &\le V_{h,\epsilon}.\label{eq:self-test-cost}
\end{align}
For the second line, the Gaussian product has sup norm at most $A^2$,
gradient norm at most $2AB$, and Laplacian norm at most $2AC+2B^2$.

\subsection{The full pair approximation theorem}

Use the same $W_X$ and $\epsilon$ for the true field and its exact
finite-source approximant, and let $R_{2,\epsilon}^{W_X}(V)$ denote
\eqref{eq:full-correlation} formed from the zeros of $V$.

\begin{theorem}[Uniform full-pair source approximation]\label{thm:full-pair}
Let $X\ge10$, $0<h\le1$, $\epsilon>0$, and suppose the complete
source error \eqref{eq:complete-delta} obeys
$\delta\le\min(1,c_0/2)$. With the explicit constants of
\eqref{eq:canonical-constants}, \eqref{eq:EX}, and
\eqref{eq:UV-costs},
\begin{equation}\label{eq:full-pair-error}
 \sup_{u,v\in\C}
 \left|R_{2,\epsilon}^{W_X}(H_0)(u,v)
          -R_{2,\epsilon}^{W_X}(\widehat H_{J,U})(u,v)\right|
 \le\frac{E_X}{\pi}
       \left(4n_XA_\epsilon U_{h,\epsilon}+V_{h,\epsilon}\right).
\end{equation}
The estimate includes every real or nonreal zero label in the window,
with its multiplicity and its actual cutoff weight.
\end{theorem}
\begin{proof}
Applying Lemma~\ref{lem:two-sided} to
$W_X(z)G_\epsilon(u-z)$ gives
\[
 \sup_u|r(u)-\widehat r(u)|\le E_XU_{h,\epsilon}/\pi.
\]
The population bound gives $0\le r,\widehat r\le2n_XA_\epsilon$.
Thus
\[
 |r(u)r(v)-\widehat r(u)\widehat r(v)|
 \le\frac{E_X}{\pi}\,4n_XA_\epsilon U_{h,\epsilon},
\]
by adding and subtracting $\widehat r(u)r(v)$. Applying the same
lemma to $W_X(z)^2G_\epsilon(u-z)G_\epsilon(v-z)$ bounds
$|d(u,v)-\widehat d(u,v)|$ by $E_XV_{h,\epsilon}/\pi$.
Subtract the self terms using Proposition~\ref{prop:weighted-pair}.
All estimates are uniform in the observation points; their Gaussian
tails do not alter the compact support of the tests in $z$.
\end{proof}

For fixed $h,X$ and $0<\epsilon\le1$, the right side is
$O_{h,X}(E_X\epsilon^{-6})$. This follows from
$A=O(\epsilon^{-2})$, $B=O(\epsilon^{-3})$,
$C=O(\epsilon^{-4})$, so $AU$ and $V$ are both
$O_h(\epsilon^{-6})$. The estimate is conservative; no optimal
width exponent is claimed. Nor does a uniform density bound imply
an $L^1$ bound on the whole noncompact observation space without
additional tail estimates.

Set $K_{X,h,\epsilon}=(4n_XA_\epsilon U_{h,\epsilon}
+V_{h,\epsilon})/\pi$. For any density tolerance $e>0$, choose
\begin{equation}\label{eq:full-precision}
 \delta_* =\min\left\{1,\frac{c_0}{2},
 M_X\left(\frac{e}{2K_{X,h,\epsilon}C_X}\right)^{1/\alpha_X}\right\}.
\end{equation}
Choose the arithmetic and radial tails in \eqref{eq:delta-H} each
at most $\delta_*/(2\mathcal B_X)$. Their explicit decay gives finite
$J,U$, and \eqref{eq:full-pair-error} is then at most $e/2$.
The remaining $e/2$ is available to a future certified evaluator.
No evaluation of the canonical pair function, or efficient runtime,
is asserted by this existence argument.

\subsection{Center averaging and the remaining source covariance}

The uniformity in observation points makes the theorem compatible
with \eqref{eq:unfolding}. If each center has an error bound
$\mathcal E_X$ from \eqref{eq:full-pair-error}, then
\begin{equation}\label{eq:averaged-transfer}
 \sup_{\zeta,\eta}
   |\overline R_2(\zeta,\eta)-\widehat{\overline R}_2(\zeta,\eta)|
       \le\int_I J_X^2\mathcal E_X\dd\pi(X).
\end{equation}
This follows by the triangle inequality. For common finite source cutoffs
at a prescribed uniform density tolerance, use the full-pair windows
\eqref{eq:full-window} with $0<h_{\min}\le h_X\le1$, together with
the positive lower bound on $\epsilon_X$ already declared. Then
$K_{X,h_X,\epsilon_X}$ and $J_X^2K_{X,h_X,\epsilon_X}$ are bounded
on the compact center interval. Finite source cutoffs can be chosen uniformly: the positive anchor is common, $M_X,\mathcal B_X,n_X,C_X$
are bounded, and $\alpha_X$ has a positive lower bound. This is a
compact-center assertion, not an all-height uniform estimate.

There is also an exact source expression for the first intensity.
Define the compact kernels
\[
 K_{X,\sigma,\zeta}(z)=
  \Delta_z\{\theta_\sigma(\Im z)W_X(z)
                   G_{\epsilon_X}(T_X\zeta-z)\}.
\]
The partition and zero-measure identity give
\begin{equation}\label{eq:source-intensity}
 \widetilde r_X(\zeta)
   =\frac{J_X}{2\pi}\sum_{\sigma=\pm}
       \int_{D(c_\sigma,r)}\log|F_{X,\sigma}(z)|
                                  K_{X,\sigma,\zeta}(z)\dd A(z).
\end{equation}
The analogous single-log expression for $\widetilde d_X$ uses
$J_X^2\Delta_z(\theta_\sigma W_X^2G_{\epsilon_X}(T_X\zeta-z)
G_{\epsilon_X}(T_X\eta-z))/(2\pi)$.
Inserting these formulas into \eqref{eq:center-covariance} identifies
a centered product of localized log integrals. Its kernels, center
law and diagonal subtraction are explicit. The individual log bounds
control approximation error; they do not establish cancellation or a
limiting shape for this center covariance. That source-specific joint
estimate is the principal next analytical question.

\section{The joint information needed for a canonical profile}\label{sec:joint-source}

Source approximation and statistical evaluation answer different questions.
The preceding theorem controls the error in replacing the entire source.
We now specify the remaining joint input and give a finite accuracy
interface for it. The derivation follows the internally assessed
fixed-center construction in \cite{OCJoint,OCJointReview}.

\subsection{Fix the law before estimating it}

For each fixed finite $T\ge10$, use the compact window
\eqref{eq:full-window} with $h=1$, and set
\begin{equation}\label{eq:fixed-regime}
 \begin{gathered}
 X\sim T^{-1}\mathbf1_{[T,2T]}\dd X,\qquad
 \ell=\min(1,4\pi/\log T),\qquad Q=[-1,1]+\ii[-1,1],\\
 \tau=\tfrac14,\qquad\epsilon=\tau\ell,\qquad J=\ell^2.
 \end{gathered}
\end{equation}
Write $r_X,d_X,R_{2,X}$ for the unfolded first intensity, self term
and factorial pair density. Thus $R_{2,X}=r_X\otimes r_X-d_X$ and
$C_T=\mathbb E R_{2,X}-(\mathbb E r_X)\otimes(\mathbb E r_X)$.
All scales are positive and fixed at each $T$; this convention alone
asserts no limiting density as $T\to\infty$.

Use common source normalization $\mathsf M=4T+26$ and put
\[
 a_* =\log(2\mathsf M/c_0),\quad
 n=\max\{1,\lceil a_*/\log(27/26)\rceil\},\quad
 m=2n,\quad A_\epsilon=(2\pi\epsilon^2)^{-1}.
\]
Here $c_0$ is the anchor of Section~\ref{sec:canonical}.
The two Jensen disks bound the number of active labels by $m$ for
both the exact and an admissible approximating field. Consequently
\begin{equation}\label{eq:joint-population}
 0\le r_X\le R:=JmA_\epsilon,\qquad
 0\le d_X\le D:=J^2mA_\epsilon^2.
\end{equation}
These size bounds do not evaluate the covariance.

There is a useful exact frequency description. For
$a_j(X)=((x_j-X)/\ell,y_j/\ell)$ define
\[
 A_X(k)=\sum_jW_X(z_j)\ee^{-\ii k\cdot a_j(X)},\qquad
 B_X(k)=\sum_jW_X(z_j)^2\ee^{-\ii k\cdot a_j(X)}.
\]
With Fourier convention $\widehat f(k)=\int\ee^{-\ii k\cdot u}f(u)\dd^2u$,
Gaussian transformation gives
\begin{equation}\label{eq:joint-fourier}
 \widehat C_T(k,q)=\ee^{-\tau^2(|k|^2+|q|^2)/2}
 \bigl[\mathbb E A_X(k)A_X(q)-\mathbb E A_X(k)\mathbb E A_X(q)
                          -\mathbb E B_X(k+q)\bigr].
\end{equation}
The covariance is bilinear, without complex conjugation. Its absolute
value is at most $m^2$ by centered Cauchy--Schwarz. Restricting both
frequency vectors to disks of radius $K$ therefore gives inverse-transform
error at most
\begin{equation}\label{eq:joint-frequency-tail}
 \frac{m^2+m}{4\pi^2\tau^4}(2\ee^{-s}-\ee^{-2s}),
 \qquad s=\tau^2K^2/2.
\end{equation}
Indeed the mass of a two-dimensional Gaussian outside such a disk is
$\ee^{-s}$ times its total mass $2\pi/\tau^2$; the complement of the
product of two disks has relative mass $2\ee^{-s}-\ee^{-2s}$.

The mixed product retains the factor
\[
 \frac1T\int_T^{2T}W_X(z_j)W_X(z_k)
             \ee^{\ii(k_x+q_x)X/\ell}\dd X
\]
and the separate marks $k_y y_j+q_y y_k$. A complex-difference projection
does not specify this general four-real-frequency transform. Translating
two labels together preserves their mutual difference, but can change
the center overlap at an endpoint of $[T,2T]$. This is an information
test on finite configurations, not another canonical zero set. The
arithmetic projection actually available is retained in
Appendix~\ref{sec:arithmetic}.

The center edges also have a quantitative cost. Replacing the center
law by its normalized restriction to $[T+1,2T-1]$ changes $C_T$ by at
most $2(3R^2+D)/T$: the deleted probability is $2/T$, the raw product
costs $R^2$ times this probability, the product of means costs $2R^2$,
and the self mean costs $D$. Deleting labels outside $[T,2T]$ gives
the same bound, as a separate comparison, since those labels can
contribute only on the edge strips. Keeping the original law incurs
neither cost.

\subsection{Regularize the logarithm, keeping the nonharmonic part}

Translate the reflected disks to $D_\sigma=D(\sigma4\ii,r)$,
$r=25/6$, $\sigma\in\{+,-\}$, and let
$f_{X,\sigma}(\xi)=F_\sigma(X+\xi)/\mathsf M$. Set
\begin{align}
 \mathcal K_{\sigma,u}(\xi)
  &=\frac{J}{2\pi}\Delta_\xi
       [\theta_\sigma W_0G_\epsilon(\ell u-\xi)],\label{eq:joint-kernel}\\
 \mathcal D_{\sigma,u,v}(\xi)
  &=\frac{J^2}{2\pi}\Delta_\xi
       [\theta_\sigma W_0^2G_\epsilon(\ell u-\xi)
                                     G_\epsilon(\ell v-\xi)].\nonumber
\end{align}
They have compact support strictly inside their disks, are independent
of $X$, and contract $\log|f_{X,\sigma}|$ to $r_X$ and $d_X$.

\begin{proposition}[Scope of harmonic replacement]\label{prop:harmonic-joint}
If a reference field is harmonic on a neighborhood of each closed
disk $D_\sigma$, every contraction of that field with the corresponding
kernel in \eqref{eq:joint-kernel} vanishes. This holds pathwise for a
random harmonic reference, including all covariance cross terms.
\end{proposition}
\begin{proof}
Twice integrating by parts moves the compact test's Laplacian onto
the reference field. Its Laplacian is zero. Thus subtracting the
reference leaves the entire contracted statistic in the residual.
\end{proof}
This proposition concerns harmonic replacements on these full planar
supports. It does not exclude contour access, nonharmonic approximations
or estimates retaining the residual. In particular the real part of
a finite holomorphic prime polynomial has zero contraction; its leading
law cannot by itself determine this planar zero statistic.

For an integer $M\ge1$ define
\[
 P_M(t)=-\frac12\sum_{a=1}^M\frac{(1-t)^a}{a},\qquad
 Y_{X,\sigma}=1-|f_{X,\sigma}|^2\in[0,1],\qquad
 L_{X,\sigma}^{(M)}=P_M(|f_{X,\sigma}|^2).
\]
The sublevel estimate \eqref{eq:sublevel}, after squaring its threshold,
has the form
$|\{|f|^2<u\}|\le\Lambda u^{1/n}$, where
\[
 \Lambda=n\pi(26/3)^2\ee^{102a_*/n},\qquad \beta=n\Lambda/2.
\]
It extends to $u\ge1$ since $\Lambda$ exceeds the disk area.
Layer cake yields
$\int\ee^{-a|f|^2}\le\Lambda\Gamma(1+1/n)a^{-1/n}$.
The gamma factor is at most one by concavity of $x^{1/n}$ under
the unit exponential probability measure. Since $(1-t)^a\le\ee^{-at}$,
the nonnegative logarithm-series remainder satisfies
\begin{equation}\label{eq:polynomial-log-error}
 \int_{D_\sigma}|L_{X,\sigma}^{(M)}-\log|f_{X,\sigma}||\dd^2\xi
 \le\frac\Lambda2\sum_{a>M}a^{-1-1/n}
 \le\beta M^{-1/n}.
\end{equation}
Zeros cause no pointwise product estimate here; their contribution is
handled by integrability and the sublevel bound. The approximation is
nonharmonic: direct differentiation gives
$\Delta P_M(|f|^2)=2M|f'|^2(1-|f|^2)^{M-1}\ge0$.

Abbreviate $U=U_{1,\epsilon}$, $V=V_{1,\epsilon}$ from
\eqref{eq:UV-costs}. Contracting the two disk errors gives
\[
 \|r_X-r_X^{(M)}\|_\infty\le d_M:=\frac{JU\beta}{\pi}M^{-1/n},
 \qquad
 \|d_X-d_X^{(M)}\|_\infty\le e_M:=\frac{J^2V\beta}{\pi}M^{-1/n}.
\]
Each product of first intensities costs $2Rd_M+d_M^2$. Hence
\begin{equation}\label{eq:joint-bias}
 \|C_T-C_T^{(M)}\|_\infty\le4Rd_M+2d_M^2+e_M,
 \qquad
 \|\overline R_2-\overline R_2^{(M)}\|_\infty
       \le2Rd_M+d_M^2+e_M.
\end{equation}

\subsection{A finite mixed-moment accuracy theorem}

Define the actual same-center moments
\[
 \mu_a^\sigma(\xi)=\mathbb E Y_{X,\sigma}(\xi)^a,\qquad
 \mu_{ab}^{\sigma\upsilon}(\xi,\omega)
   =\mathbb E Y_{X,\sigma}(\xi)^aY_{X,\upsilon}(\omega)^b.
\]
Finite expansion gives the exact surrogate
\begin{align}
 C_T^{(M)}(u,v)
  ={}&\frac14\sum_{a,b=1}^M\frac1{ab}\sum_{\sigma,\upsilon}
     \iint\mathcal K_{\sigma,u}(\xi)\mathcal K_{\upsilon,v}(\omega)
       [\mu_{ab}^{\sigma\upsilon}-\mu_a^\sigma\mu_b^\upsilon]
                                      \dd^2\xi\dd^2\omega\nonumber\\
   &+\frac12\sum_{a=1}^M\frac1a\sum_\sigma
        \int\mathcal D_{\sigma,u,v}(\xi)\mu_a^\sigma(\xi)\dd^2\xi.
                                                        \label{eq:joint-moments}
\end{align}
The plus sign on the last line comes from subtracting the negative
polynomial logarithm in the self term; the kernel is signed. The
direct full-pair formula replaces the bracket by $\mu_{ab}^{\sigma\upsilon}$.

\begin{theorem}[Finite joint-source interface]\label{thm:joint-interface}
Put $H_M=\sum_{a=1}^M a^{-1}$, $\kappa=Jr^2U$ and $\nu=J^2r^2V$.
Suppose all first moments and mixed moments in
\eqref{eq:joint-moments} are supplied with uniform errors
$\delta_1,\delta_2$, and clip estimated first moments to $[0,1]$.
The additional connected-profile error is at most
\begin{equation}\label{eq:moment-error}
 \frac{\kappa^2H_M^2}{4}(\delta_2+2\delta_1)
                 +\frac{\nu H_M}{2}\delta_1.
\end{equation}
For the direct full pair the first parenthesis can be replaced by
$\delta_2$. For a complete-source approximation with common logarithm
error $E$, the connected source error is at most
\begin{equation}\label{eq:joint-source-error}
 \frac{J^2E}{\pi}(8nA_\epsilon U+V).
\end{equation}
Adding \eqref{eq:joint-source-error}, the first bound in
\eqref{eq:joint-bias}, and \eqref{eq:moment-error} gives a valid
total connected-profile error.
\end{theorem}
\begin{proof}
The sums of the two kernels' $L^1$ norms are bounded by $\kappa$ and
$\nu$: each disk has area $\pi r^2$ and the factors in
\eqref{eq:joint-kernel} are $J/(2\pi)$ and $J^2/(2\pi)$.
Products of clipped means change by at most $2\delta_1$.
Absolute integration of the finite sums proves \eqref{eq:moment-error}.
For \eqref{eq:joint-source-error}, the unfolded first-intensity error
is $JEU/\pi$ and both intensities are bounded by $2nJA_\epsilon$.
The raw product and the product of means each cost
$4nJ^2A_\epsilon EU/\pi$; the self term costs $J^2EV/\pi$.
\end{proof}

At fixed $T$, the complete arithmetic and radial tails can make $E$
arbitrarily small, and then $M$ can make \eqref{eq:joint-bias}
arbitrarily small. The remaining input is the same-center moment array,
or a stronger estimate for its contracted combination, at the tolerance
in \eqref{eq:moment-error}. Separate one-point bounds do not supply it.
Neither a practical complexity bound nor positivity of the finite-$M$
surrogate is asserted. This theorem specifies what an acquisition must
pay; it does not report an acquired canonical profile.

\section{Newman-time regularity through collisions}\label{sec:newman-time}

Gaussian observation time and Newman time play different roles.
The former smooths a fixed divisor; the latter moves the divisor by
$\partial_\lambda H=-H_{zz}$. We retain the compact physical window
and fixed observation regime \eqref{eq:fixed-regime} throughout
$0\le\lambda\le1/2$. The result in this section is the full marked-pair
time modulus of \cite{OCTime,OCTimeReview}; it supplies conditional
transport of initial data, without assuming a positive minimum root gap.

\subsection{Complete-force inputs and cancellation of the pair singularity}

We use the canonical strip $|y_j|\le1$ and the following coarse input:
there is a finite absolute $A\ge1$ such that, throughout this interval,
\begin{equation}\label{eq:time-input}
 \left|\frac{H_\lambda'}{H_\lambda}(u+4\ii)\right|
      \le A\log(2+|u|),\qquad
 \#\{j:u\le x_j\le u+1\}\le A\log(2+|u|).
\end{equation}
The input is obtained from the Polymath asymptotic estimates and the
endpoint argument recorded in \cite{PolymathHeat2019,ParentModulus}.
At time zero the safe-line bound follows by locally uniform convergence
at each fixed $u$; the constant is uniform in $u$ before the limit.
For counts, enclose a unit interval by a slightly enlarged zero-free
rectangle, use the argument principle for small positive times, and
absorb the fixed enlargement into $A$. Reflection covers negative
heights and compactness covers the bounded-height remainder. No
numerical value of $A$ is asserted here.

For $X\ge10$ let $L_X=\log(2+X)$,
$\mathcal C_X=\{j:|x_j-X|\le3\}$, and let $m_X$ count labels in
$|x_j-X|\le1$. The input gives
\[
 n_X:=|\mathcal C_X|\le12AL_X,\qquad m_X\le2AL_X.
\]
Write $F_\lambda=H_\lambda'/H_\lambda$ and $w=x+4\ii$.
The complete exterior force at an active point $z=x+\ii y$ is
\begin{equation}\label{eq:exterior-force}
 E_X(z)=F_\lambda(w)-\sum_{k\in\mathcal C_X}\frac1{w-z_k}
  +\sum_{k\notin\mathcal C_X}
       \left(\frac1{z-z_k}-\frac1{w-z_k}\right).
\end{equation}
This retains the opposite-paired canonical logarithmic derivative.
The last series is absolutely convergent: its summands are at most
$5/|x-x_k|^2$, while the counts grow logarithmically. Unit shells
starting at distance two bound this series by $25AL_X$; the first
two terms cost at most $2AL_X$ and $n_X/3\le4AL_X$. Thus
\begin{equation}\label{eq:exterior-bound}
 |E_X(z)|\le31AL_X\qquad(|x-X|\le1,\ |y|\le1).
\end{equation}
For example the shell bound uses
$10A\sum_{h\ge2}[L_X+\log(h+2)]h^{-2}$;
$\sum_{h\ge2}h^{-2}\le1$ and the decreasing-function integral bound
\[
 \sum_{h\ge2}\frac{\log(h+2)}{h^2}
 \le\int_1^\infty\frac{\log(t+2)}{t^2}\dd t
 =\tfrac32\log3<3<\tfrac32L_X
\]
give the displayed constant. These estimates are only upper
allowances, not signed exterior information.

Let $f$ be a real $C^2$ test supported in $|x-X|\le1$, and put
$F_1=\sup_{|y|\le1}|\nabla f|$,
$F_2=\sup_{|y|\le1}\|D^2f\|_{\rm op}$. At a simple active zero,
implicit differentiation gives
\[
 \dot z_j=2\sum_{k\in\mathcal C_X,\,k\ne j}\frac1{z_j-z_k}
                  +2E_X(z_j).
\]
For $d=x_j-x_k$, $e=y_j-y_k$, the two ordered force terms combine
into $-2K_f(z_j,z_k)$, where
\begin{align}
 K_f(z_j,z_k)
  &=\frac{e(f_y(z_j)-f_y(z_k))-d(f_x(z_j)-f_x(z_k))}{d^2+e^2}\nonumber\\
  &=\int_0^1\frac{e^2 f_{yy}-d^2 f_{xx}}{d^2+e^2}
                     (z_k+t(z_j-z_k))\dd t.\label{eq:chord-cancellation}
\end{align}
The mixed derivatives cancel, and the chord stays in the strip.
Hence $|K_f|\le F_2$ even as the separation tends to zero.
Only pairs with at least one active label contribute, at most
$m_Xn_X$ unordered pairs. The exterior costs $2m_XF_1\,31AL_X$.
We have proved, at simple times,
\begin{equation}\label{eq:compact-time-test}
 \left|\partial_\lambda\sum_j f(z_j)\right|
       \le A^2L_X^2(48F_2+124F_1).
\end{equation}
This is the useful cancellation: bounding the two velocities separately
would introduce a reciprocal root separation that the paired test does
not require.

\subsection{The full pair and connected moduli}

Set $f_{X,u}(z)=JW_X(z)G_\epsilon(X+\ell u-z)$ and define
$r_{\lambda,X},d_{\lambda,X},R_{2,\lambda,X}$ as before.
On the canonical strip the transverse cutoff is identically one,
including its first two derivatives at the boundary. The horizontal
window has derivative bounds $15/4$ and $24$. The Gaussian has
gradient bound $A_\epsilon/(\epsilon\sqrt{\ee})$ and Hessian operator
bound $A_\epsilon/\epsilon^2$. Define
\begin{align}
 a_0&=\frac{J}{2\pi\epsilon^2}=\frac8\pi,&
 a_1&=a_0\left(\frac{15}4+\frac1{\epsilon\sqrt{\ee}}\right),\nonumber\\
 a_2&=a_0\left(24+\frac{15}{2\epsilon\sqrt{\ee}}+\frac1{\epsilon^2}\right),&
 L&=\log(2+2T),\quad R_*=2ALa_0,\nonumber\\
 L_1&=A^2L^2(48a_2+124a_1),&
 L_d&=A^2L^2[96(a_0a_2+a_1^2)+248a_0a_1],\nonumber\\
 L_2&=2R_*L_1+L_d,&L_C&=4R_*L_1+L_d.
                                                        \label{eq:time-constants}
\end{align}

\begin{theorem}[Full marked-pair time modulus]\label{thm:time-modulus}
For $\lambda,\mu\in[0,1/2]$, uniformly in $X\in[T,2T]$,
\[
 \begin{aligned}
 \|r_{\lambda,X}-r_{\mu,X}\|_\infty&\le L_1|\lambda-\mu|,\\
 \|d_{\lambda,X}-d_{\mu,X}\|_\infty&\le L_d|\lambda-\mu|,\\
 \|R_{2,\lambda,X}-R_{2,\mu,X}\|_\infty&\le L_2|\lambda-\mu|.
 \end{aligned}
\]
The same first and pair bounds hold after center averaging, and
$\|C_\lambda-C_\mu\|_\infty\le L_C|\lambda-\mu|$.
The norms are on $Q$ or $Q^2$, and the estimates include multiple-zero times.
\end{theorem}
\begin{proof}
The product rule gives $|f_{X,u}|\le a_0$,
$|\nabla f_{X,u}|\le a_1$, $\|D^2f_{X,u}\|\le a_2$.
The self test $f_{X,u}f_{X,v}$ has gradient at most $2a_0a_1$ and
Hessian at most $2a_0a_2+2a_1^2$. Formula
\eqref{eq:compact-time-test} proves $L_1,L_d$ at simple times.
Also $r_{\lambda,X}\le R_*$. The raw product costs $2R_*L_1$;
subtracting the self term gives $L_2$. The connected statistic subtracts
a second product, now of averaged means, giving $L_C$.

It remains to justify integration through events. The complete family
is analytic in $(\lambda,z)$ and is not identically zero on a slice,
since $H_\lambda(0)>0$. In the padded compact rectangle
$[T-2,2T+2]+\ii[-2,2]$, local zero-free boundaries and Rouch\'e's theorem
give continuous finite multisets of roots, with multiplicity preserved.
The weighted first, self and pair sums are therefore jointly continuous.

At a zero of multiplicity $q$ at $(\lambda_0,z_0)$, put
$t=\lambda-\lambda_0$, $\zeta=z-z_0$. The heat equation makes the
lowest parabolic-degree Taylor polynomial a nonzero multiple of
\[
 \ee^{-t\partial_\zeta^2}\zeta^q
  =(2t)^{q/2}\operatorname{He}_q(\zeta/\sqrt{2t}).
\]
The polynomial identity follows by expanding the finite differential
series, or the Hermite generating function \cite{DLMFHermiteOC}.
The $q$ Hermite zeros are real and distinct: the Rodrigues expression
$\operatorname{He}_q(x)=(-1)^q\ee^{x^2/2}(\dd/\dd x)^q\ee^{-x^2/2}$
and integration by parts give orthogonality to every polynomial of
degree below $q$ for the positive Gaussian weight. If there were fewer
than $q$ real sign changes, their linear factors would form such a
polynomial whose product with $\operatorname{He}_q$ has constant sign,
contradicting orthogonality. Rescaling and applying Rouch\'e
near each one shows that all nearby roots are simple for sufficiently
small nonzero $t$; their displacements are
$\sqrt{2t}\,h_j+O(|t|)$. Thus multiple events are isolated in this
compact space-time set and there are finitely many. Integrate the
uniform simple-time estimates on the intervening intervals and use
continuity at their endpoints. The interval lengths sum to
$|\lambda-\mu|$, so no event-count factor appears. Averaging over
the unchanged center law preserves the estimates.
\end{proof}

\subsection{Exterior truncation and conditional transport}

Truncate only the absolutely convergent difference series in
\eqref{eq:exterior-force}, retaining terms with $|x-x_k|<M$, for
integer $M\ge2$. Unit counts and a decreasing-function integral give
\begin{equation}\label{eq:exterior-tail}
 e_M(X)=\frac{10A}{M-1}[L_X+\log(M+1)+1].
\end{equation}
Indeed the shell majorant is
$10A[L_X+\log(h+2)]/h^2$; its integral from $M-1$ to infinity
is bounded by \eqref{eq:exterior-tail}, using
$\log(1+s)\le s$ after rescaling. The induced scalar derivative
error is at most $2m_XF_1e_M(X)$.
Let $m_*=2AL$, and replace $L_X$ by $L$ in $e_M$ to define $e_M^*$.
The full-pair and connected derivative errors are at most, respectively,
\begin{align}
 B_{2,M}&=(4R_*m_*a_1+4m_*a_0a_1)e_M^*,\nonumber\\
 B_{C,M}&=(8R_*m_*a_1+4m_*a_0a_1)e_M^*.
                                                        \label{eq:pair-exterior-error}
\end{align}
The first coefficient pays two first-intensity derivative errors;
the second pays the self-product derivative. Connected subtraction
doubles the first coefficient. Uniform convergence allows time
integration, giving errors $B_{2,M}|\lambda-\mu|$ and
$B_{C,M}|\lambda-\mu|$. Root locations and undifferentiated intensities
remain those of the exact divisor; this is not an error theorem for
an independently truncated root dynamics.

If a function $F_0$ approximates $\overline R_{2,\mu}$ within $e_0$,
then
\begin{equation}\label{eq:initial-transport}
 \|\overline R_{2,\lambda}-F_0\|_\infty
       \le e_0+L_2|\lambda-\mu|.
\end{equation}
The connected version uses $L_C$. For large $T$, the displayed allowances
are $L_1,L_d=O_A((\log T)^4)$ and $L_2,L_C=O_A((\log T)^5)$.
They supply no initial profile and no favorable sign.

Lipschitz continuity is compatible with a kink at collision. For
$p_\lambda(z)=z^2-2\lambda$ and constant local window, the ordered
two-label density at observations $u=v=0$ is exactly
\[
 2(2\pi\epsilon^2)^{-2}\ee^{-2|\lambda|/\epsilon^2}.
\]
Its right-minus-left slope is
$-8(2\pi\epsilon^2)^{-2}/\epsilon^2$. At $\lambda=0$ both labels
remain and the factorial coefficient is $2$, not zero. This local
heat control proves neither a canonical collision nor differentiability
through one; it shows why the theorem claims a modulus rather than a
globally smooth closed equation.

\section{Gaussian height windows and an explicit marked projection}\label{sec:gaussian-window}

We now change the physical window, while retaining the observation
Gaussian and the uniform center law. This change makes center integration
explicit and is useful for finite acquisition. Its tails must be paid;
the compact-window time constants of Section~\ref{sec:newman-time}
do not automatically apply to it. The analytical construction is from
\cite{OCGaussianWindow,OCProjection}.

\subsection{An exact center integral with controlled tails}

For $h>0$ use
\begin{equation}\label{eq:gaussian-height-window}
 W_X^{\rm G}(x+\ii y)=\ee^{-(x-X)^2/(2h^2)}\kappa_0(y),
 \qquad X\sim\operatorname{Unif}[T,2T].
\end{equation}
Here $h$ is a standard deviation, not the support radius of the earlier
compact window. For $m=(x+x')/2$ and $d=x-x'$,
\begin{equation}\label{eq:height-center-integral}
 \begin{gathered}
 \mathbb E_X\ee^{-(x-X)^2/(2h^2)}\ee^{-(x'-X)^2/(2h^2)}
   =I_{T,h}(m)\ee^{-d^2/(4h^2)},\\
 I_{T,h}(m)=\frac{h\sqrt\pi}{2T}
 \left[\operatorname{erf}\frac{2T-m}{h}
                    -\operatorname{erf}\frac{T-m}{h}\right].
 \end{gathered}
\end{equation}
The identity follows from
$(x-X)^2+(x'-X)^2=2(m-X)^2+d^2/2$ and a one-dimensional Gaussian integral.
Both finite center edges remain in $I_{T,h}$.

The full pair density has a similar exact center integration before
any projection. Put
$\sigma^2=h^2\epsilon^2/(h^2+\epsilon^2)$,
$\alpha=h^2/(h^2+\epsilon^2)$ and
$a_u(z)=x-\alpha\ell\Re u$. Completing the square gives
\[
 JW_X^{\rm G}(z)G_\epsilon(X+\ell u-z)
  =p_u(z)\ee^{-(X-a_u(z))^2/(2\sigma^2)},
\]
where
\[
 p_u(z)=\frac{J\kappa_0(y)}{2\pi\epsilon^2}
  \exp\left[-\frac{(y-\ell\Im u)^2}{2\epsilon^2}
             -\frac{\ell^2(\Re u)^2}{2(h^2+\epsilon^2)}\right].
\]
For each ordered distinct-label pair, multiply $p_u(z)p_v(w)$ by
the right side of \eqref{eq:height-center-integral} with $h=\sigma$,
$m=(a_u(z)+a_v(w))/2$, and $d=a_u(z)-a_v(w)$.
This simplifies integration over a supplied divisor; it does not
estimate the unknown divisor or its joint source field.

From \eqref{eq:time-input}, the total height weight is bounded by
$B_0(X)$, and the weight at $|x_j-X|\ge R$ by $B_R(X)$, where
\begin{equation}\label{eq:gaussian-weight-tail}
 B_R(X)=2A\sum_{n=R}^\infty
       [\log(2+X)+\log(n+2)]\ee^{-n^2/(2h^2)},\qquad R\in\mathbb Z_{\ge0}.
\end{equation}
Indeed each shell is covered by two unit intervals and their logarithms
are at most $\log(2+X)+\log(n+2)$. These bounds are uniform in
$\lambda\in[0,1/2]$. With $a_0=J/(2\pi\epsilon^2)$, keeping only
the fixed slab $[T-R,2T+R]$ loses at most
\begin{equation}\label{eq:gaussian-pair-tail}
 2a_0^2B_0(2T)B_R(2T)\quad\hbox{in the full pair},\qquad
 4a_0^2B_0(2T)B_R(2T)\quad\hbox{in the connected pair}.
\end{equation}
Every omitted pair has at least one omitted weight; connected subtraction
pays the product of means a second time.

\subsection{What linear Gaussian smoothing does to the source}

For the normalized real Gaussian $\gamma_h$ of variance $h^2$,
absolute integration of the complete cosine representation gives
\begin{equation}\label{eq:source-heat-shift}
 \int_\R\gamma_h(X-x)H_\lambda(x+\ii y)\dd x
       =H_{\lambda-h^2/2}(X+\ii y).
\end{equation}
Each frequency is multiplied by $\ee^{-h^2u^2/2}$.
This is a linear identity. It does not pass through the logarithm.
For $p_\lambda(z)=z^2-2\lambda$ at $\lambda=h^2/2$, the convolved
polynomial is $z^2$, whereas the original zeros are $\pm h$.
At the origin the Gaussian mean of $\log|x^2-h^2|$ is finite and
the log modulus of the convolved polynomial is $-\infty$.

The log-source formulas themselves remain valid for the noncompact
height window. To justify this, let
\[
 M_0=\max\left\{1,\int_0^\infty\ee^{u^2/2+11u}\Phi(u)\dd u\right\}.
\]
Then $|H_\lambda|\le M_0$ on every disk $D(q+4\ii,7)$, and
integration of \eqref{eq:time-input} along the safe line gives
$\log|H_\lambda(q+4\ii)|\ge\log H_0(4\ii)-A|q|\log(2+|q|)$.
Subharmonic circular means, integrated in radius, therefore imply
\begin{equation}\label{eq:noncompact-log-tail}
 \int_{[q,q+1]\times[-2,2]}|\log|H_\lambda||\dd^2z
 \le49\pi[2\log M_0-\log H_0(4\ii)+A|q|\log(2+|q|)].
\end{equation}
The rectangle fits inside the disk. The pointwise bound
$|\log|H||\le2\log M_0-\log|H|$ and the mean lower bound prove
the inequality, including the negative log near zeros.

The Laplacians of $f_{X,u}=JW_X^{\rm G}G_\epsilon(X+\ell u-\cdot)$
and $f_{X,u}f_{X,v}$ have Gaussian horizontal decay times a polynomial
of degree at most two; their transverse support stays in $[-2,2]$.
Multiply them first by a horizontal cutoff with uniformly bounded first
two derivatives. Formula \eqref{eq:noncompact-log-tail} makes the
cutoff error summable and vanishing, while
\eqref{eq:gaussian-weight-tail} handles the zero side. Passing to the
limit proves the absolutely convergent formulas
\begin{equation}\label{eq:gaussian-log-contraction}
 r_X(u)=\frac1{2\pi}\int\log|H_\lambda|\Delta f_{X,u},\qquad
 d_X(u,v)=\frac1{2\pi}\int\log|H_\lambda|\Delta(f_{X,u}f_{X,v}).
\end{equation}
Uniform absolute bounds justify products and the center average.
Their joint input is still $\mathbb E_X[L_X(\xi)L_X(\omega)]$,
with $L_X(\xi)=\log|H_\lambda(X+\xi)|$, or its covariance.
A whole-line comparison also needs the approximating field's log tail.
Neither \eqref{eq:source-heat-shift} nor the tail bounds acquire a
sharp estimate of this joint field.

\subsection{Integrate one marked observation test exactly}

Return to $\lambda=0$, keep $\epsilon/\ell=1/4$, and use the Gaussian
height window. The following bounded test preserves two separate
imaginary marks:
\[
 \Psi(u,v)=\exp\left[-\frac{(\Re u-\Re v)^2}2
                     -\frac{(\Im u-1/2)^2}2
                     -\frac{(\Im v+1/2)^2}2\right].
\]
Define the projection over the \emph{entire} observation plane by
\begin{equation}\label{eq:marked-projection}
 P_T=\mathbb E_X\int_{\C^2}\Psi(u,v)R_{2,X}(u,v)\dd^2u\dd^2v.
\end{equation}
The Jacobian $J=\ell^2$ makes each unweighted observation Gaussian
have mass one. A real difference has variance $2(1/4)^2=1/8$ and
each transverse observation has variance $1/16$. Using
\[
 \mathbb E\ee^{-(a+Z)^2/2}
   =(1+v)^{-1/2}\ee^{-a^2/[2(1+v)]},\qquad Z\sim N(0,v),
\]
which follows by completing the square, the three convolution factors
multiply to $c_G=32\sqrt2/51$. Thus
\begin{equation}\label{eq:projection-kernel-general}
 P_T=\sum_{j\ne k}K_{T,h,\ell}(z_j,z_k),\qquad
 K_{T,h,\ell}(z,w)=c_G I_{T,h}(m)
   \ee^{-[1/(4h^2)+4/(9\ell^2)]d^2}b_+(y)b_-(y'),
\end{equation}
where $b_\pm(y)=\exp[-\frac8{17}(y/\ell\mp1/2)^2]$.
Here $\kappa_0=1$ on the canonical strip. Tonelli's theorem applies
to the nonnegative sums and integrals, with finiteness supplied by
\eqref{eq:gaussian-weight-tail}.

For the finite acquisition below set $h=\ell=1$, $\epsilon=1/4$,
and write $I_T=I_{T,1}$, $K=K_{T,1,1}$. Then
\begin{equation}\label{eq:projection-kernel}
 K(z,w)=c_G I_T(m)\ee^{-25d^2/36}b_+(y)b_-(y'),\qquad
 D(z):=K(z,z)=c_G I_T(x)\ee^{-4/17}\ee^{-16y^2/17}.
\end{equation}
The ordering matters:
$K(z,w)/K(w,z)=\exp[\frac{16}{17}(y-y')]$.
The midpoint dependence and the two marks therefore cannot be replaced
by a function of the complex difference alone. This is a specified
scalar consumer of the full hierarchy, not a reconstruction of it.

\section{Finite acquisition: centroid cancellation and cell ranges}\label{sec:acquisition}

The preceding projection admits two rigorous finite acquisition methods.
The first uses counts and first moments, and is insensitive to internal
root gaps. The second uses only counts and certified cell geometry, and
can exploit Gaussian decay when cells are separated. Their inputs differ;
neither method is uniformly superior. We derive both error bounds before
giving the finite census as supporting evidence in
Appendix~\ref{app:status}. The source arguments are
\cite{OCCentroids,OCCellRanges}.

\subsection{A complete theta tail with explicit constants}

We first pay for every zero outside the retained rectangle, without
requiring a numerical value of the coarse constant $A$.
The source majorant \eqref{eq:source-basic-majorant} gives
$\Phi(u)<10\ee^u r_u^2\ee^{-r_u}$: consecutive terms in
$\sum_{l\ge1}l^4\ee^{-\pi(l^2-1)}$ have ratio at most
$16\ee^{-3\pi}<1/2$.
Since $u^2\le r_u/12$, substitution $v=r_u$ gives
\begin{equation}\label{eq:explicit-source-majorant}
 \mathcal M:=\int_0^\infty\ee^{u^2}\Phi(u)\dd u
 \le\frac{10}{4\pi^{1/4}}\int_\pi^\infty v^{5/4}\ee^{-11v/12}\dd v
 <5(12/11)^3<7.
\end{equation}
The inequality $u^2\le r_u/12$ follows from $\ee^{2u}\ge2u$ and
$\pi>3$; for the last estimate use $v^{5/4}\le v^2$ and extend the
integral to zero. On $0\le u\le1/10$, $3<r_u<8$, and the first
positive source term is at least $9\ee^{-8}$. Thus
$H_0(0)>1/7290$ and $\log(\mathcal M/H_0(0))<20$.

For $|z|\le2r$, $|\cos(zu)|\le\ee^{2ru}\le\ee^{r^2+u^2}$.
Jensen's formula on the disk of radius $2r$ now gives
\begin{equation}\label{eq:global-count-jensen}
 N_{\rm disk}(r)\log2\le r^2+20,
 \qquad N_{\rm disk}(r)\le2r^2+40.
\end{equation}
Boundary zeros follow by a larger-radius limit. Each real unit slab
starting at $q$ lies, on the canonical strip, in a disk of radius
$|q|+2$; it therefore contains at most $16(|q|+2)^2$ labels.

For $T=1000$ retain
$\mathcal R=(988,2012)+\ii(-3/2,3/2)$.
An omitted label has $|x_j-X|\ge12$ at every center.
Two unit slabs per shell bound its total Gaussian height weight by
\begin{equation}\label{eq:explicit-weight-tail}
 b=32\sum_{n=12}^\infty(2003+n)^2\ee^{-n^2/2}<2^{-44}.
\end{equation}
To verify the strict bound, write $n=12+j$ and dominate the summand by
$2015^2(1+j)^2\,2^{-72}2^{-12j}$. For $q=2^{-12}$,
$\sum_{j\ge0}(1+j)^2q^j=(1+q)/(1-q)^3<2$ and $2015<2^{11}$.
If the complete retained count is $N\le1024$, then
\begin{equation}\label{eq:explicit-projection-tail}
 0\le P_{1000}-P_{1000}^{\mathcal R}
       \le2c_G(N+b)b<2^{-32}.
\end{equation}
This includes negative heights and both infinite horizontal tails.
The count cap is a stated input; the Gaussian decay does not certify
the retained census by itself.

\subsection{Counts and first moments give a second-order error}

On $\mathcal R$, $Z(z)=\zeta(1/2-\ii z/2)$ has the same divisor as
$H_0$ and has no pole or trivial zero. Its completion factors are
holomorphic and nonzero there. On a zero-free cell boundary,
the argument principle and its first-moment version give
\begin{equation}\label{eq:contour-centroid}
 n_C=\frac1{2\pi\ii}\oint_{\partial C}\frac{Z'}Z\dd z,
 \qquad q_C=\frac1{2\pi\ii}\oint_{\partial C}(z-a_C)\frac{Z'}Z\dd z,
 \qquad c_C=a_C+q_C/n_C.
\end{equation}
Residues show that $n_C$ counts labels and $c_C$ is their exact
unweighted complex centroid. A zero-free completion multiplier
contributes zero to both contour integrals. These are classical
contour identities; the useful issue is the uncertainty they remove
from the specified pair projection.

\begin{lemma}[Kernel derivative bounds]\label{lem:kernel-derivatives}
For the smooth real extension of \eqref{eq:projection-kernel} to
$\C^2$, with real two-coordinate derivatives,
\[
 \|\nabla_zK\|,\|\nabla_wK\|\le9/T,\quad
 \|D_z^2K\|_{\rm op},\|D_w^2K\|_{\rm op}\le18/T,
 \quad\|\nabla D\|\le4/T,\quad\|D^2D\|_{\rm op}\le10/T.
\]
\end{lemma}
\begin{proof}
Before center integration the kernel is
\[
 \frac{c_G}{T}\int_T^{2T}w(x-X)w(x'-X)g(x-x')b_+(y)b_-(y')\dd X,
 \quad w(s)=\ee^{-s^2/2},\quad g(d)=\ee^{-4d^2/9}.
\]
These four Gaussian factors and their first two derivatives have
absolute value at most one. In $z$ derivatives, keep $w(x'-X)$
and integrate it over the whole line, with mass $\sqrt{2\pi}<3$.
Gradient components are bounded by $6/T,3/T$, and Hessian entries
in $xx,xy,yy$ by $12/T,6/T,3/T$. The Euclidean gradient and
maximum absolute row sum give $9/T,18/T$. The other variable is identical.
For $D$, use $|I_T|,|I_T'|\le2/T$, $|I_T''|\le8/T$,
and derivative bounds $1,2$ for $\ee^{-16y^2/17}$.
The Hessian entries are bounded by $8/T,2/T,4/T$, proving the claim.
\end{proof}

\begin{theorem}[Centroid acquisition error]\label{thm:centroid}
Partition the retained divisor into convex cells with counts $n_C$,
centroids $c_C$ and total count $N$. Set
$S_2=\sum_C\sum_{z_j\in C}|z_j-c_C|^2$ and
\[
 P_{\rm bar}=\sum_{C,E}n_Cn_EK(c_C,c_E)-\sum_Cn_CD(c_C).
\]
Then
\begin{equation}\label{eq:centroid-error}
 |P_T^{\mathcal R}-P_{\rm bar}|\le\frac{18N+5}{T}S_2.
\end{equation}
In particular cells of diameter at most $r$ give
$S_2\le Nr^2$. If $|\widehat c_C-c_C|\le\sigma$, replacing
centroids by these approximations costs at most
$(18N^2+4N)\sigma/T$.
\end{theorem}
\begin{proof}
For a real $C^2$ test $F$, Taylor expansion in a cell gives
\[
 \left|\sum_{z_j\in C}F(z_j)-n_CF(c_C)\right|
 \le\tfrac12\sup_C\|D^2F\|_{\rm op}
                         \sum_{z_j\in C}|z_j-c_C|^2,
\]
since $\sum_{z_j\in C}(z_j-c_C)=0$ cancels both real components
of the linear term. Apply this successively to both arguments of the
raw all-label sum; each application costs $9NS_2/T$. The self term
costs $5S_2/T$. Convexity puts the centroid inside its cell, so
$S_2\le Nr^2$. Gradient bounds applied to both arguments and then
the diagonal prove the rounding allowance.
\end{proof}
No simple-root or gap hypothesis enters this theorem. In a collapsed
cell the factorial coefficient is $n_C(n_C-1)$, and the cost of
collapsing distinct roots is already paid by \eqref{eq:centroid-error}.
The weights stay inside $K$; replacing these unweighted centroids
by an unexplained weighted average would lose the cancellation.

\subsection{Direct ranges exploit separation and Gaussian decay}

Suppose instead that each of $N$ rectangular cells contains exactly one
label, and that they are sorted horizontally with adjacent interval
distance at least $s>0$. Let $u$ and $v$ bound their horizontal and
vertical widths. For $C<E$ form the exact intervals
\[
 M_{CE}=[(x_C^-+x_E^-)/2,(x_C^++x_E^+)/2]=[m_-,m_+],
 \qquad D_{CE}=[x_E^--x_C^+,x_E^+-x_C^-]=[d_-,d_+].
\]
Here $d_->0$. Define ideal, exact-real endpoints
\begin{align}
 A_-&=\frac{\sqrt\pi}{2T}
   [\operatorname{erf}(2T-m_+)-\operatorname{erf}(T-m_-)],\nonumber\\
 A_+&=\frac{\sqrt\pi}{2T}
   [\operatorname{erf}(2T-m_-)-\operatorname{erf}(T-m_+)],\nonumber\\
 I_-&=\max(0,A_-),\quad I_+=\min(2/T,A_+),\qquad
 G_-=\ee^{-25d_+^2/36},\quad G_+=\ee^{-25d_-^2/36}.
                                                        \label{eq:cell-endpoints}
\end{align}
Monotonicity of $\operatorname{erf}$ gives the enclosure of $I_T$;
clipping uses $0\le I_T\le\sqrt\pi/T<2/T$.
For each mark let $s_{C,\pm}^-,s_{C,\pm}^+$ be the minimum and maximum
of $(y\mp1/2)^2$ on the cell's vertical interval, and put
\[
 B_{C,\pm}^-=\ee^{-8s_{C,\pm}^+/17},\qquad
 B_{C,\pm}^+=\ee^{-8s_{C,\pm}^-/17},\qquad
 H_{CE}^\pm=B_{C,+}^\pm B_{E,-}^\pm+B_{E,+}^\pm B_{C,-}^\pm.
\]
The last expression sums both ordered mark products. Set
\begin{equation}\label{eq:ideal-cell-sum}
 L_*=c_G\sum_{C<E}I_-G_-H_{CE}^-,\qquad
 U_*=c_G\sum_{C<E}I_+G_+H_{CE}^+.
\end{equation}
All factors are nonnegative, so $P_T^{\mathcal R}\in[L_*,U_*]$.

\begin{theorem}[Gaussian range-width bound]\label{thm:cell-ranges}
Under these singleton-cell and separation hypotheses, the ideal radius obeys
\begin{equation}\label{eq:range-width}
 \frac{U_*-L_*}{2}
   \le\frac{N(10u+4v)}{T}\sum_{k=1}^{N-1}\ee^{-(ks)^2/4}
   <\frac{N(10u+4v)\sqrt\pi}{sT}
\end{equation}
when $N\ge2$ and $u+v>0$. For $N\le1$ the retained factorial sum is zero.
\end{theorem}
\begin{proof}
At index separation $k$, every admissible horizontal gap is at least
$ks$. Put $q_k=\ee^{-(ks)^2/4}$. Endpoint differentiation gives
$\operatorname{wid}I\le2u/T$ and $\sup I\le2/T$.
For $g(d)=\ee^{-25d^2/36}$,
\[
 g(d)\le\ee^{-d^2/4},\qquad |g'(d)|\le2\ee^{-d^2/4}.
\]
Indeed write $a=25/36$, use
$|g'|=2\sqrt a(\sqrt a|d|\ee^{-ad^2/2})\ee^{-ad^2/2}$,
$\sup_{t\ge0}t\ee^{-t^2/2}=\ee^{-1/2}$ and $a/2>1/4$.
Since the gap interval has width at most $2u$,
$\sup G\le q_k$ and $\operatorname{wid}G\le4u q_k$.
Each mark is at most one and has derivative at most
$\sqrt{16/(17\ee)}<1$, so its interval width is at most $v$.
Telescoping the product of nonnegative intervals shows that the width
of one ordered contribution is at most
\[
 \frac{2u}{T}q_k+\frac2T\,4u q_k+\frac2Tq_k(2v)
       =\frac{10u+4v}{T}q_k.
\]
There are $2(N-k)$ ordered pairs of this index separation. Summing
and dividing by two proves the first bound. The decreasing Gaussian
has sum strictly below its integral from zero to infinity,
$\sqrt\pi/s$, proving the second.
\end{proof}
The arbitrary-$s$ form makes explicit the geometric input to the
specialized estimate in \cite{OCCellRanges}. It is a sufficient error
bound for separated singleton cells, not a uniform canonical gap theorem.
For general counts, cross-cell coefficients become $n_Cn_E$;
same-cell coefficients $n_C(n_C-1)$ require an enclosure of $K$ on
$C\times C$, not only the diagonal $D$. Such cells do not inherit
\eqref{eq:range-width} without a new width analysis.

\begin{corollary}[Separate arithmetic and geometric allowances]\label{cor:range-arithmetic}
Suppose enclosing balls for the exact endpoints $L_*,U_*$ have radii
$r_L,r_U$, and let $a$ be the lower endpoint of the first ball and $b$
the upper endpoint of the second. If the omitted tail is at most $t$,
then $[a-t,b+t]$ encloses the complete projection and has radius at most
\begin{equation}\label{eq:range-total}
 \frac{U_*-L_*}{2}+r_L+r_U+t.
\end{equation}
\end{corollary}
\begin{proof}
A ball containing an exact endpoint can place its outward edge as much
as twice its radius away from that endpoint. Thus
$b-a\le U_*-L_*+2r_U+2r_L$. The symmetric tail enlargement adds $t$
to the radius. The nonnegative tail would also allow the sharper
asymmetric enlargement $[a,b+t]$; we retain the conservative formula.
\end{proof}
Endpoint arithmetic, the certified census, and the analytic tail are
three distinct premises. A width theorem does not evaluate either
endpoint. The supporting example below retains this distinction.

\section{Scope, controls, and outstanding tasks}\label{sec:outstanding}

The main results now form a quantitative chain. The factorial construction
specifies the complete observable and its self correction. Source
stability bounds the error in replacing it. The mixed-moment theorem
specifies the further joint information needed for a canonical profile.
The Newman-time theorem transports supplied initial information through
collisions. Finally, exact Gaussian integration turns one declared
projection into a finite acquisition problem with separate spatial,
geometric and arithmetic allowances. None of these steps supplies the
missing canonical joint estimate by itself.

The contribution is measured by the displayed error bounds, not by the
number of representations or formal identities. The standard factorial,
contour and Gaussian identities are used as tools; no literature novelty
is claimed for them. Negative results retain their stated classes:
harmonic replacements on the full planar supports, linear Gaussian
convolution passed through a logarithm, and the particular information
lost by a difference-only projection.

The collision controls catalogued in the repository were screened
analytically against these claims; their earlier numerical
certificates were not rerun for this manuscript.
Generic identities and approximation estimates are expected to
hold for controls satisfying their hypotheses. The missing
canonical premise in the following table explains why a control
does not become a theta example; it does not establish a
favorable theta estimate.

\begin{center}\small
\begin{tabularx}{\textwidth}{@{}>{\raggedright\arraybackslash}p{0.29\textwidth}Y@{}}
\toprule
Control family & Premise absent from a canonical inference\\
\midrule
Euler-factor modifications &
The unchanged complete theta coefficients and conductor. Preserved
real zeros do not account for added nonreal zeros.\\
Weighted-square sources and finite coefficient or tail changes &
The complete normalized reflection and unchanged arithmetic
source. Finite positivity or endpoint agreement is insufficient.\\
Reflected super-Gaussians &
The integer-square theta arithmetic, despite positivity and
reflection symmetry.\\
Smooth shape and positive-primitive controls &
The full canonical coefficient expansion, despite matching
shape or auxiliary-transform properties.\\
Finite-prime orbit controls &
The full reflection and all-coefficient source identity;
finite prime data do not determine the residual divisor.\\
Atomic event and finite-profile controls, with continuous extensions &
The full theta profile and coefficient structure; finite
observations do not certify the complete integer zero divisor.\\
RH-conditional normalized endpoint controls &
Ordinary Dirichlet-series arithmetic. The full correlation retains
added double zeros with multiplicity; no simplicity premise is used.\\
Variable observation kernels &
The cosine oscillator and the corresponding Newman equation;
fixed amplitude does not restore the altered observation law.\\
Actual-theta three-profile relaxation &
The coupling of all observations to an integer zero divisor;
the relaxation permits fractional masses.\\
\bottomrule
\end{tabularx}
\end{center}

The polynomial controls in
\eqref{eq:pair-control} and \eqref{eq:real-divisor-control}
are additional local information tests, not members of the
unchanged theta family. The catalog and the project's scope
rules are cited in the accompanying manuscript guide and
continuation note, with source paths rather than new claims
of universal obstruction.

The outstanding tasks are the following, in order of analytical importance.

\begin{enumerate}[leftmargin=*,itemsep=0.55em]
\item \textbf{Acquire genuine same-center canonical information.}
For the fixed compact regime \eqref{eq:fixed-regime}, obtain the moments
in \eqref{eq:joint-moments}, or a sharper estimate of their contracted
combination, with the explicit error in Theorem~\ref{thm:joint-interface}.
For the Gaussian regime, the analogous target is the joint field in
\eqref{eq:gaussian-log-contraction}. A one-point estimate or another
identity for the covariance does not satisfy this task. A paid full
profile is a worthwhile result independently of a sine law or an RH
application.

\item \textbf{Make the time allowance informative for supplied data.}
Theorem~\ref{thm:time-modulus} is complete at fixed compact window and
positive blur. A practical transport application still needs an initial
profile, a quantitative value or usable bound for $A$, and a time interval
on which the additive allowance is useful. Extending this theorem to a
Gaussian physical window requires summing the moving-zero and exterior
costs with its tails. It is not obtained merely by substituting the
Gaussian window into the compact constants. Time-dependent windows or
unfolding scales also require their own derivative terms.

\item \textbf{Extend quantitative transfer to higher orders and limits.}
The hierarchy is defined at every finite order, but higher canonical
error bounds must subtract every repeated-label partition and control
all Gaussian product derivatives. Growing order requires a simultaneous
order-tail bound. Removing blur returns a measure, not a pointwise
smooth density; growing height changes the source and time constants.
Normalization by factorial mass or first intensities requires positive
lower bounds and errors for those denominators.

\item \textbf{Complete the declared finite projection as supporting evidence.}
The inherited census and Theorem~\ref{thm:cell-ranges} pay the geometric
allowance for $P_{1000}$. The endpoint evaluation still has to return
outward enclosures with $r_L+r_U\le1/1000$ and include the complete
$2^{-32}$ tail. The fixed E13 r2 contract permits that bounded evaluation;
this manuscript supplies no scalar value. A failure of that arithmetic
gate would be an execution result, not a failure of the analytical
width theorem or of all centroid methods.

\item \textbf{Supply an additional consumer before drawing collision conclusions.}
A regular full profile can coexist with a collision, as the quadratic
heat control shows. Any collision application needs a resolving
observation scale and a proved lower signal, together with an upper or
signed comparison that beats it. Conditioning, an initial gap, or a
signed exterior estimate must be stated and paid where used. The
collision-constraint investigation owns that separate priority.

\item \textbf{Assess the integrated claims at their actual evidence levels.}
The E07, E08 and E12 acceptances retain their original scopes.
The Gaussian-window and acquisition arguments, the arbitrary-separation
form of \eqref{eq:range-width}, and this assembly have proof-author
checking only. A separate claim audit should check their interfaces
and constants. The earlier B02 and B04 assessments E03 and E04 remain
pending; successful typesetting and finite bookkeeping do not replace them.
\end{enumerate}

The appendices preserve the useful comparison results and their limits,
then record the finite evidence and its outstanding arithmetic gate.

\appendix
\section{An ensemble benchmark derived by marginalization}\label{app:ensemble}

To distinguish a model joint law from the canonical construction,
consider the circular density
\begin{equation}\label{eq:cue-joint}
 p_N^{\rm circ}(\boldsymbol\theta)
   =\frac1{N!(2\pi)^N}
      \prod_{j<k}|\ee^{\ii\theta_j}-\ee^{\ii\theta_k}|^2,
 \qquad 0\le\theta_j<2\pi.
\end{equation}
This is a familiar determinant-product example
\cite[Section 4]{Borodin2009DPP}; its relevant calculation is short.
Put $\phi_m(\theta)=\ee^{\ii m\theta}/\sqrt{2\pi}$ for $0\le m<N$.
The Vandermonde identity writes \eqref{eq:cue-joint} as
$|\det[\phi_m(\theta_j)]|^2/N!$. On expanding both determinants,
orthonormality kills all unequal permutations after full integration,
so the density has mass one. If only $N-k$ coordinates are integrated,
orthogonality instead retains the $k\times k$ minors; expanding their
determinants gives
$\rho_k=\det[K_N(\theta_a,\theta_b)]_{a,b\le k}$, where
\[
 K_N(\theta,\phi)=\frac1{2\pi}\sum_{m=0}^{N-1}
                                      \ee^{\ii m(\theta-\phi)}.
\]
In particular $\rho_2=(N/(2\pi))^2-|K_N|^2$.
Unfolding by $x=N\theta/(2\pi)$ and summing the geometric series yields
\begin{equation}\label{eq:cue-sine}
 \rho_{2,N}^{\rm unf}(s)
    =1-\left(\frac{\sin\pi s}{N\sin(\pi s/N)}\right)^2
       \longrightarrow 1-\left(\frac{\sin\pi s}{\pi s}\right)^2
\end{equation}
for fixed $s$, with removable values interpreted by continuity.
If each coordinate of the limiting real-line pair density is smoothed
by $g_\epsilon$, its separation profile becomes
\begin{equation}\label{eq:smoothed-sine-reference}
 1-g_{\sqrt2\epsilon}*\left(\frac{\sin\pi s}{\pi s}\right)^2.
\end{equation}
This follows from the same difference-variance calculation as
Proposition~\ref{prop:separation}. These equations provide a reference
for what a model marginal calculation can produce. Neither the
Vandermonde law nor its sine profile has been derived for the localized
canonical zero configuration in this manuscript.

\section{Arithmetic pair information and its localization limits}\label{sec:arithmetic}

The full smoothed pair function retains a common center, a separation,
and two transverse marks. We now compare that information with one known
arithmetic projection. The comparison explains which variables an existing
theorem integrates out, and therefore which new estimates the present
program needs. It also supports the marked scalar specialization in
Appendix~\ref{sec:pairs}. The conclusions concern the stated imports;
they are not a limitation theorem for all arithmetic correlation methods.

\subsection{The corrected cumulative and dyadic imports}

Write $L=\log T$, $B_T=TL/(2\pi)$, and use $r_0>0$ for the
exponential base in the arithmetic formula. For an ordinate interval
$I$ define
\begin{equation}\label{eq:arithmetic-F}
 \mathcal F_I(r_0)=
 \sum_{\gamma_j,\gamma_k\in I}
       r_0^{\rho_j-\rho_k}W(\rho_j-\rho_k),
 \qquad W(u)=\frac4{4-u^2}.
\end{equation}
All zeta zeros and all multiplicity copies are included, along
with their self labels. The corrected source is
Baluyot--Goldston--Suriajaya--Turnage-Butterbaugh
\cite[Section 3]{BGST2026PairI}, version 3.
The first line of its cumulative formula (3.5) gives, uniformly
for $1\le r_0\le T$,
\begin{equation}\label{eq:bgst-cumulative-raw}
 \mathcal F_{(0,T]}(r_0)
 =\frac{T}{2\pi r_0^2}\{L^2+O(L)\}
       +\frac{T}{2\pi}\log r_0+O(T\sqrt L).
\end{equation}
Consequently, for $0\le\alpha\le1$,
\begin{equation}\label{eq:bgst-cumulative}
 \frac{\mathcal F_{(0,T]}(T^\alpha)}{B_T}
 =T^{-2\alpha}\{L+O(1)\}+\alpha+O(L^{-1/2}).
\end{equation}
The corresponding dyadic theorem, its equation (3.2), is
\begin{equation}\label{eq:bgst-dyadic}
 \frac{\mathcal F_{(T,2T]}(T^\alpha)}{B_T}
 =T^{-2\alpha}L+\alpha
       +O(T^{-2\alpha}\sqrt L+L^{-1/2}).
\end{equation}
The stronger first error in \eqref{eq:bgst-cumulative} must not
be transferred to the different population in
\eqref{eq:bgst-dyadic}. These are imported arithmetic estimates,
not reproved prime mean-square theorems in this draft. Their
uniform big-$O$ constants and onsets are not numerically
instantiated here.

For clarity, positivity of the full sum has an elementary finite
explanation. If $u_j=\rho_j-1/2$, the zero multiset in $I$ is
invariant under $u\mapsto-\overline u$. The identity
$W(u)=\int_\R\ee^{-2|t|}\ee^{tu}\dd t$ for $|\Re u|<2$
therefore gives
\begin{equation}\label{eq:arithmetic-gram}
 \mathcal F_I(r_0)=
 \int_\R\ee^{-2|t|}
       \left|\sum_{\gamma_j\in I}
                   \ee^{(\log r_0+t)u_j}\right|^2\dd t\ge0.
\end{equation}
Interchange is valid for the finite population and
$|\Re(\rho_j-\rho_k)|\le1$.
Swapping labels shows
$\mathcal F_I(r_0^{-1})=\mathcal F_I(r_0)$.
Neither observation makes every complex summand nonnegative.

\subsection{A triangular-frequency projection and its actual value}

Let $q_{\rm freq}(\alpha)=(1-|\alpha|)_+$ and define its entire
transform
\begin{equation}\label{eq:fejer-transform}
 K_T(u)=\int_{-1}^1(1-|\alpha|)\ee^{\alpha Lu}\dd\alpha
       =\left(\frac{\sinh(Lu/2)}{Lu/2}\right)^2,\qquad K_T(0)=1.
\end{equation}
One can verify the formula by integrating twice, or by observing
that the triangular function is the convolution of two unit
interval indicators and multiplying their exponential transforms.
The ordered factorial statistic for the dyadic population
$\mathcal Z(T)=\{\rho:T<\gamma\le2T\}$ is
\begin{equation}\label{eq:A_T}
 A_T=\frac1{B_T}\sum_{\substack{\rho_j,\rho_k\in\mathcal Z(T)\\j\ne k}}
       W(\rho_j-\rho_k)K_T(\rho_j-\rho_k).
\end{equation}
Writing $n_T=|\mathcal Z(T)|$, the removed self mass is exactly
$n_T/B_T$. Distinct labels at a multiple zero remain.

\begin{proposition}[Value of the imported projection]\label{prop:bulk-value}
The dyadic import \eqref{eq:bgst-dyadic} and the classical
unconditional count $n_T=B_T+O(T)$ imply
\begin{equation}\label{eq:bulk-value}
 A_T=\frac13+O(L^{-1/2}).
\end{equation}
If the error in \eqref{eq:bgst-dyadic} is bounded by
$C_{\rm BG}(T^{-2\alpha}\sqrt L+L^{-1/2})$ and
$|n_T/B_T-1|\le C_N/L$ in a common range, then
\begin{equation}\label{eq:bulk-error}
 \left|A_T-\frac13\right|
 \le\frac{2C_{\rm BG}}{\sqrt L}+\frac{C_N+1/2}{L}.
\end{equation}
\end{proposition}
\begin{proof}
Integrate \eqref{eq:bgst-dyadic} against the triangular weight,
use evenness, and subtract the self mass. The two leading integrals
are
\[
 2L\int_0^1(1-\alpha)\ee^{-2L\alpha}\dd\alpha
       =1-\frac{1-T^{-2}}{2L},\qquad
 2\int_0^1(1-\alpha)\alpha\dd\alpha=\frac13.
\]
The integrated first error is at most
$2C_{\rm BG}\sqrt L\int_0^1\ee^{-2L\alpha}\dd\alpha
\le C_{\rm BG}/\sqrt L$.
The second contributes $C_{\rm BG}/\sqrt L$.
The zero count and the remaining $1/(2L)$ term give
\eqref{eq:bulk-error}. In particular the frequency-zero spike
has been retained.
\end{proof}

This is an actual arithmetic value for the declared projection.
It is not the value of the marked test in Appendix~\ref{sec:pairs}.
In the $H$ coordinate, the population lies in $2T<x\le4T$,
a window of half-width $T$ centered at $3T$.
Moreover
\[
 L(\rho_j-\rho_k)=
 \frac L2\bigl(-(y_j-y_k)+\ii(x_j-x_k)\bigr),
\]
so the microscopic difference scale is $4\pi/L$ in the $x$
coordinate. The projection retains a complex difference but
neither arbitrary common-center localization nor the two
separate transverse marks.

A simple $H$-pair at $x_*\pm\ii d$, $0<d<1$, has zeta
real-part difference $d$. Its two ordered mutual terms in
\eqref{eq:A_T} contribute
\begin{equation}\label{eq:arithmetic-defect}
 \frac2{B_T}\frac4{4-d^2}
       \left(\frac{\sinh(dL/2)}{dL/2}\right)^2 .
\end{equation}
For $d=c/L$ with fixed $c>0$, this is $O((TL)^{-1})$,
while the available error allowance in \eqref{eq:bulk-error}
is of order $L^{-1/2}$ and the main value is $1/3$.
At $d=0$ two different coincident real labels contribute $2/B_T$.
Other complex summands need not be nonnegative, so
\eqref{eq:arithmetic-defect} is neither a lower bound on $A_T$
nor the total change produced by inserting a pair: cross
interactions remain.
Thus this known projection supplies neither the desired
short-window marked value nor a useful absolute single-defect
comparison with the stated error.

\subsection{What the real-root and replacement-method benchmarks supply}

The classical Montgomery pair-correlation framework
\cite{Montgomery1973} motivates distinguishing bulk difference
information from a fully marked local marginal. An elementary
real-root identity illustrates the distinction. For a polynomial
with real roots $a_1,\ldots,a_n$, counted with multiplicity, put
\[
 p_P(x)=\sum_{j=1}^n\frac{y}{\pi((x-a_j)^2+y^2)},\qquad y>0 .
\]
The Cauchy densities satisfy $P_y*P_y=P_{2y}$; this follows,
for example, from their Fourier transform $\ee^{-y|\xi|}$.
Therefore
\begin{equation}\label{eq:real-poisson}
 2\pi y\norm{p_P}_{L^2(\R)}^2-n
 =\sum_{j\ne k}\frac{4y^2}{(a_j-a_k)^2+4y^2}.
\end{equation}
Indeed the integral of the $j,k$ product is
$P_{2y}(a_j-a_k)$, and each self term becomes one after
multiplication by $2\pi y$.
At $y=h/2$ the separation kernel is exactly $c_h$.
This real-population benchmark, used in the earlier heat-flow
collision investigation, does not establish that an
approximation to the canonical field has only real roots.
Discarding its possible complex zeros would remove the very
information under investigation.

Likewise the local replacement principle of Tao and Vu
\cite[Theorems 2.1 and 3.1]{TaoVu2014LocalZeros}
is useful as an architecture, not an imported theorem about
$H_0$. Its hypotheses include nondegeneracy, non-clustering,
and joint smooth-test comparison of logarithmic values, in a
degree-at-most-$n$ polynomial setting; the real version also
requires reference-field repulsion bounds.
A compact-source entire transform is not a polynomial, and
its local Jensen count is not its polynomial degree.
Theorem~\ref{thm:pair-transfer} avoids that probabilistic
comparison problem for deterministic accuracy on one disk,
but supplies no comparison in law or favorable reference value.

The earlier Rodgers--Tao-related local heat-flow work, as
discussed in the predecessor comparison \cite{OCArithmetic},
supplies a further methodological lesson: real simple populations,
exterior forces, and same-root correlation errors must retain
their exact hypotheses and combined accounting. None of those
real-population hypotheses is produced here for a field with
possibly nonreal zeros. The separate compact-window argument in Section~\ref{sec:newman-time}
now gives a time modulus without a real-simple-population premise.
It does not supply a closed signed correlation hierarchy.

\section{Reflection and finite configuration algebra}\label{sec:configuration}

Let a finite labelled configuration be invariant under conjugation.
Suppose $|y_j|\le1$, and choose
$w_j=w((x_j-X)/h)$ with $0\le w\le1$ and compact support, $h>0$.
Only labels in that support matter. Define the horizontal moment
measures and their first useful masses by
\begin{equation}\label{eq:qmoments}
 q_\ell=\sum_jw_jy_j^\ell\delta_{x_j},\qquad
 N_w=q_0(\R),\qquad Q_w=q_2(\R).
\end{equation}
Conjugation cancels each odd moment at each horizontal coordinate,
not merely after integration:
$q_{2r+1}=0$. This elementary observation determines exactly how much
information a quadratic transverse difference can retain.

\begin{proposition}[Quadratic pair reduction]\label{prop:quadratic}
Let $\chi$ be bounded, even and nonnegative. Then
\begin{align}
 S_2(\chi)&=\sum_{j\ne k}
    w_jw_k\chi(x_j-x_k)(y_j-y_k)^2 \notag\\
 &=2\int(\chi*q_0)(x)\dd q_2(x)
 \le 2\norm{\chi}_\infty N_wQ_w .\label{eq:quadratic-reduction}
\end{align}
In particular $S_2(1)=2N_wQ_w$.
\end{proposition}
\begin{proof}
The terms with $j=k$ vanish and can be restored. Expanding the square
gives the two terms $q_2\otimes q_0$ and $q_0\otimes q_2$,
minus $2q_1\otimes q_1$, all tested against $\chi(x-x')$.
The last term is zero. Evenness of $\chi$ makes the other two equal.
Their positivity gives the upper bound.
\end{proof}

If $\chi(0)=1$ and there are $m$ copies at each of $x_*\pm\ii d$,
their two ordered cross populations contribute
$8m^2w((x_*-X)/h)^2d^2$ to $S_2$.
They contribute $2mw((x_*-X)/h)d^2$ to $Q_w$.
Thus the test detects a transverse pair, but its exact reduction
prevents interpreting it as independently acquired two-point
information.

The same mechanism extends to all even difference powers. For $k\ge1$,
\begin{equation}\label{eq:higher-moments}
 S_{2k}(\chi)=
 \sum_{r=0}^k \binom{2k}{2r}
       \iint\chi(x-x')\dd q_{2k-2r}(x)\dd q_{2r}(x').
\end{equation}
This follows by binomial expansion; all odd factors vanish.
Since $|y_j-y_k|\le2$,
$S_{2k}(\chi)\le2^{2k-2}S_2(\chi)$.
These are identities and bounds for this particular family of tests;
they do not reconstruct the full marked pair measure.

\subsection{A configuration generating function with an order bound}

Put $a_j=w_jy_j^2\in[0,1]$ and
\begin{equation}\label{eq:configuration-product}
 Z(\eta)=\prod_{j=1}^N(1+\eta a_j)
        =\sum_{k=0}^N e_k\eta^k,\qquad
 p_r=\sum_j a_j^r,\qquad p_1=Q_w.
\end{equation}
The $e_k$ are elementary symmetric polynomials of the $a_j$.
Differentiating the finite product gives, first near $\eta=0$,
\[
 \frac{Z'(\eta)}{Z(\eta)}
 =\sum_j\frac{a_j}{1+\eta a_j}
 =\sum_{r\ge1}(-1)^{r-1}p_r\eta^{r-1}.
\]
Comparing coefficients proves Newton's recurrence
\begin{equation}\label{eq:newton}
 e_0=1,\qquad
 ke_k=\sum_{r=1}^k(-1)^{r-1}p_re_{k-r}.
\end{equation}
It is exact finite algebra; no closure assumption enters.

\begin{proposition}[Product comparison and growing-order truncation]\label{prop:product}
For $\eta\ge0$,
\begin{equation}\label{eq:product-bounds}
 \log(1+\eta)Q_w\le\log Z(\eta)\le\eta Q_w .
\end{equation}
For every integer $K\ge0$,
\begin{equation}\label{eq:order-tail}
 0\le Z(\eta)-\sum_{k=0}^Ke_k\eta^k
 \le \ee^{\eta Q_w}
       \frac{(\eta Q_w)^{K+1}}{(K+1)!},
\end{equation}
where $e_k=0$ for $k>N$. If $Q_w\le B$ and $0<\varepsilon<1$,
the right side is at most $\varepsilon$ whenever $k=K+1$ satisfies
\begin{equation}\label{eq:order-choice}
 k\ge\max\left\{1,\ 2\ee\eta B,\
           \frac{\eta B+\log(1/\varepsilon)}{\log2}\right\}.
\end{equation}
\end{proposition}
\begin{proof}
On $a\in[0,1]$, concavity of $\log(1+\eta a)$ puts its graph
above the chord $a\log(1+\eta)$ and below its tangent $\eta a$.
Summing proves \eqref{eq:product-bounds}.
Expansion of $(\sum a_j)^k$ contains every product with distinct labels
$k!$ times and additional nonnegative terms, so $e_k\le Q_w^k/k!$.
For $t\ge0$,
$\sum_{j=k}^\infty t^j/j!\le t^k\ee^t/k!$, since
$(k+\ell)!\ge k!\ell!$. This proves \eqref{eq:order-tail}.
Finally $k!\ge(k/\ee)^k$ makes its right side at most
$\ee^{\eta B}2^{-k}$ under the first two size conditions, and the
last condition gives $\varepsilon$.
\end{proof}

For complex $\eta$, use the logarithm branch equal to zero at $\eta=0$.
If $q=|\eta|\max_j a_j<1$, the same finite product gives
\begin{equation}\label{eq:complex-log-tail}
 \left|\log Z(\eta)-
       \sum_{r=1}^K\frac{(-1)^{r-1}}r\,p_r\eta^r\right|
 \le
 \frac{|\eta|^{K+1}p_{K+1}}{(K+1)(1-q)} .
\end{equation}
For each factor, expand $\log(1+\eta a_j)$, use
$1/r\le1/(K+1)$ in the tail, and sum the geometric series.
Thus increasing the factorial order is permitted only with a
simultaneous radius or tail payment, not from fixed-order identities
alone.

A local count can supply $B$. For example, if the population in
every relevant interval of length at most one is bounded by
$A_{\mathrm{count}}\log(3+|X|+h)$, and
$\supp w\subset[-1,1]$, covering $[X-h,X+h]$ gives
\[
 Q_w\le N_w\le
 A_{\mathrm{count}}(\lceil2h\rceil+1)\log(3+|X|+h).
\]
This is a conditional count corollary, with its constant retained.
The canonical disk estimate in Section~\ref{sec:canonical} supplies a
different explicit population bound in its stated regime.

\subsection{Defect signals and the cost of averaging}

On a plateau where $w=1$, a simple pair $x_*\pm\ii d$ contributes
\begin{equation}\label{eq:scalar-signals}
 Q_w\ge2d^2,\qquad
 \log Z(\eta)\ge2\log(1+\eta d^2).
\end{equation}
All other terms are nonnegative. Both observables are zero for every
all-real configuration, including real multiple zeros. They test
nonreality, not simplicity.

If the center $X$ is averaged uniformly on an interval of length $L$
which contains the full translate of the window around $x_*$, the
pair's respective contributions become
\begin{align}
 &\frac{2d^2h}{L}\int_\R w(t)\dd t,\qquad
 \frac{2h}{L}\int_\R\log(1+\eta d^2w(t))\dd t,\label{eq:average-scalar}\\
 &\frac{8d^2h}{L}\int_\R w(t)^2\dd t
       \quad\text{for its mutual contribution to }S_2
       \text{ when }\chi(0)=1 .\notag
\end{align}
These follow by the substitution $t=(x_*-X)/h$.
Averaged and every-center statements therefore have different
single-pair margins. With complete height and transverse-scale
coverage, vanishing of these nonnegative tests is simply a
reformulation of the assertion that all canonical zeros are real,
equivalently RH. That reformulation supplies no new estimate and is
not supported by sampling checks.

\section{Scalar acquisition as a supporting specialization}\label{app:scalar}

Choose $0<h\le1$, a nonnegative $C^2$ cutoff $w$, supported in
$[-1,1]$ and equal to one on $[-1/2,1/2]$, with $w\le1$.
Choose $0\le\kappa\le1$, equal to one on $[0,1]$ and zero on
$[2,\infty)$, with bounded first and second derivatives. Define
\begin{equation}\label{eq:scalar-tests}
 a(x,y)=w((x-X)/h)y_+^2\kappa(y),\qquad
 \varphi_Q=2a,\qquad
 \varphi_\eta=2\log(1+\eta a)\quad(\eta\ge0).
\end{equation}
These compact $C^{1,1}$ tests are supported in
$|x-X|\le1$, $0\le y\le2$. The support lies strictly within
$D(c,r)$ since $\sqrt{17}<25/6$. At $y=0$ the value and the first
normal derivative vanish, so their weak Laplacians contain no
boundary delta.

The canonical strip, conjugation, and \eqref{eq:zero-measure} imply
\begin{equation}\label{eq:scalar-fold}
 Q_{X,h}=\nu_F(\varphi_Q),\qquad
 \log Z_{X,h}(\eta)=\nu_F(\varphi_\eta).
\end{equation}
The factor two accounts for lower conjugate copies. Real labels
contribute zero. The harmonic logarithm of $D$ disappears.
Define the comparison values
$\widehat Q=\nu_{\widehat F}(\varphi_Q)$ and
$\widehat L_\eta=\nu_{\widehat F}(\varphi_\eta)$.
The approximating zeros in $1<y<2$ keep their actual cutoff weights:
there is no strip theorem for the approximant.

\begin{theorem}[Complete one-point acquisition]\label{thm:scalar}
For $X\ge10$, $0<h\le1$, and the tests \eqref{eq:scalar-tests},
choose $J,U$ so that \eqref{eq:complete-delta} is at most
$\min(1,c_0/2)$. Then
\begin{equation}\label{eq:scalar-transfer}
 |Q_{X,h}-\widehat Q|\le K_QE_X,\qquad
 |\log Z_{X,h}(\eta)-\widehat L_\eta|\le K_\eta E_X,
 \qquad
 K_\varphi=\frac{\norm{\Delta\varphi}_\infty}{2\pi}.
\end{equation}
All multiplicities are included. The constants are explicitly
given by \eqref{eq:radii}, \eqref{eq:canonical-constants} and
\eqref{eq:complete-delta}.
\end{theorem}
\begin{proof}
The anchor and disk bounds establish Theorem~\ref{thm:log};
apply Corollary~\ref{cor:weak-transfer} and \eqref{eq:scalar-fold}.
\end{proof}

For $v(y)=y_+^2\kappa(y)$ the derivative costs are directly available:
\begin{align}
 \Delta a&=h^{-2}w''v+wv'',&
 |\nabla a|^2&=h^{-2}|w'|^2v^2+w^2|v'|^2,\label{eq:scalar-derivatives}\\
 \Delta\varphi_\eta&=
 2\left(\frac{\eta\Delta a}{1+\eta a}
              -\frac{\eta^2|\nabla a|^2}{(1+\eta a)^2}\right).
 \notag
\end{align}
In particular $K_Q\le(h^{-2}\norm{w''}_\infty\norm v_\infty+
\norm w_\infty\norm{v''}_\infty)/\pi$.
The denominators in \eqref{eq:scalar-derivatives} are at least one;
growing $\eta$ still incurs its displayed derivative cost.
The Jensen count also gives the coarse canonical bound
$N_w\le2n_X$, hence $Q_w\le2n_X$, because every upper label in
the window lies in $D(c,\rho_*)$, every lower nonreal label has an
upper partner, and real labels need be counted only once.

For a nonzero test, any desired statistic accuracy $e>0$ is achieved
by setting
\begin{equation}\label{eq:precision-choice}
 \delta_*=\min\left\{1,\frac{c_0}2,\
    M_X\left(\frac{e}{2K_\varphi C_X}\right)^{1/\alpha_X}\right\}
\end{equation}
and choosing finite cutoffs satisfying
\begin{equation}\label{eq:tail-allocation}
 5C_2(J)I_{b_*}\le\frac{\delta_*}{2\mathcal B_X},
 \qquad B_0I_{b_*}^{>U}\le\frac{\delta_*}{2\mathcal B_X}.
\end{equation}
Then $K_\varphi E_X\le e/2$. This proves finite precision exists,
without claiming a uniform computation time. If an evaluator returns
$\widehat q$ with $|\widehat q-\widehat Q|\le E_{\rm eval}$,
the actual exclusion condition for a plateau pair of height at least
$d$ is
\begin{equation}\label{eq:scalar-consumer}
 \widehat q+E_{\rm eval}+K_QE_X<2d^2 .
\end{equation}
There is an analogous condition with margin
$2\log(1+\eta d^2)$. Quadrature and rounding belong in
$E_{\rm eval}$. No favorable $\widehat q$ is obtained from the
approximation theorem itself.

\section{A separation-dependent marked pair statistic}\label{sec:pairs}

The reduction in Proposition~\ref{prop:quadratic} motivates retaining
both transverse marks separately. The following test is nonnegative,
vanishes on real configurations, and includes a nonconstant separation
kernel. Its approximation follows from the same one-variable logarithm
estimate, without introducing a probability law.

\subsection{A fixed test and its exact response}

Define the piecewise polynomial cutoffs
\begin{equation}\label{eq:p-cutoff}
 p(t)=
 \begin{cases}
 1,&t\le0,\\
 1-10t^3+15t^4-6t^5,&0<t<1,\\
 0,&t\ge1,
 \end{cases}
 \qquad w(t)=p(2|t|-1),\qquad v(t)=t_+^2p(t-1).
\end{equation}
Thus $w=1$ on $[-1/2,1/2]$, $w=0$ off $[-1,1]$,
and $v(t)=t^2$ for $0\le t\le1$, with $v=0$ off $[0,2]$.
The function $v$ is $C^{1,1}$; smoothness of higher order at zero
is neither asserted nor needed.

For $X\ge10$ and $0<h,b\le1$, let
\begin{align}
 a(z)&=w((x-X)/h)v(y/b),&
 q(z)&=w((x-X)/h)v(|y|/b),\label{eq:pair-marks}\\
 c_h(u)&=\frac{h^2}{u^2+h^2},&
 K(z,z')&=a(z)a(z')c_h(x-x').\notag
\end{align}
For any conjugation-symmetric holomorphic field defined on a
neighborhood of the support of $q$, set
\begin{equation}\label{eq:marked-statistic}
 S_{X,h,b}(H)=\sum_{j\ne k}q(z_j)q(z_k)c_h(x_j-x_k).
\end{equation}
All multiplicity copies are separate labels. The center is fixed;
there is no implicit averaging over $X$.

\begin{proposition}[Marked-pair signal]\label{prop:pair-signal}
Every term in \eqref{eq:marked-statistic} is nonnegative.
A simple conjugate pair $x_*\pm\ii d$ with
$|x_*-X|\le h/2$ and $0<d\le b$ forces
\begin{equation}\label{eq:marked-signal}
 S_{X,h,b}(H)\ge2(d/b)^4 .
\end{equation}
For $m$ copies at each conjugate location, their population contributes
$(4m^2-2m)(d/b)^4$. Every all-real configuration gives zero,
regardless of multiplicity.
\end{proposition}
\begin{proof}
On the plateau each of the $2m$ labels has weight $(d/b)^2$.
Every difference of their horizontal coordinates is zero, so
$c_h=1$. There are $2m(2m-1)$ ordered distinct-label pairs.
Other summands cannot reduce their contribution. Real labels
have $q=0$.
\end{proof}

A useful declared scale is
\begin{equation}\label{eq:matched-scale}
 h=b=h_X:=\min(1,4\pi/\log X).
\end{equation}
This uses the leading mean-spacing scale in the $x=2\gamma$
coordinate; it is a definition, not an exact spacing assertion.
The band $b/2\le d\le b$ has signal at least $1/8$.
A single such band does not cover arbitrarily small transverse
displacements or all displacements up to the canonical strip edge.

\subsection{Conjugate folding and the self-index correction}

The support of $a$ is contained in $D(c,r)$ with the same buffer as
in Appendix~\ref{app:scalar}. For a nonzero holomorphic $V$ on
the disk, write $\mu_V=\nu_V$ there and define
\[
 T_V(K)=\iint K\dd\mu_V\dd\mu_V,\qquad
 D_V(a^2)=\int a^2\dd\mu_V .
\]
If $V$ has the same zeros on this disk as a conjugation-symmetric
full field, then
\begin{equation}\label{eq:folded-pair}
 S(V)=4T_V(K)-2D_V(a^2).
\end{equation}
Each upper nonreal label represents itself and its lower conjugate,
so the full tensor square has factor four. Removing the self label
of each of these two copies has factor two. At an upper point of
multiplicity $m$, the tensor has mass $m^2$ and the one-point
correction removes $m$, rather than $m^2$. In particular the two
conjugate partners are different full-plane labels even when their
upper representatives coincide.

Applying \eqref{eq:weak-onepoint} separately in the two independent
variables gives
\begin{align}
 T_V(K)&=\frac1{4\pi^2}
  \iint\log|V(z)|\log|V(z')|\,
       \Delta_z\Delta_{z'}K(z,z')\dd A\dd A',\label{eq:pair-log}\\
 D_V(a^2)&=\frac1{2\pi}
       \int\log|V|\,\Delta(a^2)\dd A .\notag
\end{align}
For this particular $K$ the mixed weak derivative is bounded:
each factor $a$ is differentiated at most twice. Product logarithms
are integrable by Fubini, including near the geometric diagonal,
because they belong to different variables. Mollification in the
support buffer proves the formula for the stated cutoffs.
Adding a harmonic term to either logarithm changes neither integral;
integrate first in its corresponding variable. Thus scalar
normalizations and zero-free multipliers are immaterial.

\subsection{Tensor transfer with only second derivatives}

\begin{theorem}[Deterministic pair transfer]\label{thm:pair-transfer}
Let $f,g$ satisfy the normalized disk and anchor hypotheses of
Theorem~\ref{thm:log}, with the disk translated to $c$, and suppose
\[
 \int_{D(c,r)}|\log|f|-\log|g||\dd A\le E .
\]
Let $K$ be a real symmetric $C_c^{1,1}$ test on $D(c,r)^2$.
Its sections and diagonal trace have bounded weak Laplacians.
Choose bounds
\[
 D_2\ge\max\{\norm{\Delta_zK}_\infty,
                  \norm{\Delta_{z'}K}_\infty\},\qquad
 D_{\rm diag}\ge\norm{\Delta[K(z,z)]}_\infty .
\]
For the folded functional
$\mathcal S_V(K)=4T_V(K)-2\int K(z,z)\dd\mu_V(z)$,
\begin{equation}\label{eq:pair-transfer}
 |\mathcal S_f(K)-\mathcal S_g(K)|
 \le\frac{E}{\pi}\bigl(4nD_2+D_{\rm diag}\bigr),
\end{equation}
where $n$ is \eqref{eq:local-constants}. For the test
\eqref{eq:pair-marks} and conjugation-symmetric full fields,
this is the difference of their statistics
\eqref{eq:marked-statistic}.
\end{theorem}
\begin{proof}
Each relevant population is at most $n$ by the Jensen part of
Theorem~\ref{thm:log}. For fixed $z'$ in the relevant projection
of the support, the weak one-variable identity gives
\[
 |(\mu_f-\mu_g)(K(\,\cdot\,,z'))|
 \le\frac{ED_2}{2\pi}.
\]
The exact telescoping identity
\[
 \mu_f\otimes\mu_f-\mu_g\otimes\mu_g
 =(\mu_f-\mu_g)\otimes\mu_f+
                         \mu_g\otimes(\mu_f-\mu_g)
\]
allows us to integrate this bound against positive zero measures.
The two masses cost at most $2n$, giving
$|T_f-T_g|\le nED_2/\pi$.
The diagonal trace is a one-variable test, so its error is at most
$ED_{\rm diag}/(2\pi)$. Multiplication by four and two proves
\eqref{eq:pair-transfer}.
\end{proof}

The population bound replaces a fourth-derivative cost. Indeed if
$D_4\ge\norm{\Delta_z\Delta_{z'}K}_\infty$ and
$M=\pi r^2A$, \eqref{eq:log-mass} and
\[
 L_f(z)L_f(z')-L_g(z)L_g(z')
 =(L_f-L_g)(z)L_f(z')+L_g(z)(L_f-L_g)(z')
\]
give the valid alternative
\begin{equation}\label{eq:pair-fourth}
 |\mathcal S_f-\mathcal S_g|
 \le E\left(\frac{2MD_4}{\pi^2}
                      +\frac{D_{\rm diag}}{\pi}\right).
\end{equation}
Theorem~\ref{thm:pair-transfer} instead uses only second derivatives.
This compares two explicit allowances, not optimal approximation
rates. Strong deterministic accuracy on the same disk suffices
for this transfer; an independently favorable pair value remains
a different question.

\subsection{Explicit constants and the canonical theorem}

Differentiating \eqref{eq:p-cutoff},
$p'=-30t^2(1-t)^2$ and
$p''=-60t(1-t)(1-2t)$ on $(0,1)$.
Their maxima satisfy
$\norm{p'}_\infty\le15/8$ and
$\norm{p''}_\infty\le10/\sqrt3<6$.
Consequently
\[
 \norm w_\infty\le1,\quad \norm{w'}_\infty\le15/4,\quad
 \norm{w''}_\infty\le24,\qquad
 \norm v_\infty\le4,\quad \norm{v'}_\infty\le23/2,\quad
 \norm{v''}_\infty\le41.
\]
For example on $1<t<2$,
$v'=2tp+t^2p'$ and
$v''=2p+4tp'+t^2p''$, which give the last two bounds.
Thus
\begin{equation}\label{eq:a-bounds}
 |a|\le A_0=4,\quad |a_x|\le A_x=15/h,\quad
 |a_y|\le A_y=23/(2b),\quad
 |\Delta a|\le D_a=96h^{-2}+41b^{-2}.
\end{equation}
The partial fraction formula
$c_h(u)=\frac{h}{2\ii}[(u-\ii h)^{-1}-(u+\ii h)^{-1}]$
implies $|c_h^{(j)}(u)|\le j!h^{-j}$ on the real axis.
Since
\[
 \Delta_zK=(\Delta a)a'c_h+2a_xa'c_h'+aa'c_h'',
 \qquad K(z,z)=a(z)^2,
\]
where $a'=a(z')$ in this display only, admissible constants are
\begin{align}
 D_2&=A_0D_a+2A_xA_0/h+2A_0^2/h^2
       =536h^{-2}+164b^{-2},\label{eq:D2}\\
 D_{\rm diag}
   &=2A_0D_a+2A_x^2+2A_y^2
       =1218h^{-2}+\frac{1185}{2}b^{-2}.\label{eq:Ddiag}
\end{align}
For \eqref{eq:pair-fourth}, expansion of the mixed Laplacian permits
the choice
\[
 D_4= D_a^2+\frac{4A_xD_a}{h}
 +\frac{4A_0D_a+8A_x^2}{h^2}
 +\frac{24A_xA_0}{h^3}+\frac{24A_0^2}{h^4}
 =D_a^2+76D_a/h^2+3624/h^4.
\]
No fourth derivative of $v$ has been used.

\begin{theorem}[Complete canonical marked-pair acquisition]\label{thm:canonical-pair}
Let $X\ge10$, $0<h,b\le1$, and define the test by
\eqref{eq:p-cutoff}--\eqref{eq:marked-statistic}.
Choose finite $J,U$ so that the complete error
\eqref{eq:complete-delta} satisfies \eqref{eq:delta-requirement}.
Then
\begin{align}
 |S_{X,h,b}(H_0)-S_{X,h,b}(\widehat H_{J,U})|
 &\le E_{\rm pair},\label{eq:canonical-pair}\\
 E_{\rm pair}
 &=\frac{C_X(\delta/M_X)^{\alpha_X}}{\pi}
 \left[
 4n_X(536h^{-2}+164b^{-2})
 +1218h^{-2}+\frac{1185}{2}b^{-2}
 \right].\notag
\end{align}
In particular, when $b=h$,
\begin{equation}\label{eq:matched-error}
 \boxed{\displaystyle
 E_{\rm pair}=\frac{2800n_X+3621/2}{\pi h^2}\,
                  C_X(\delta/M_X)^{2/(n_X+2)} .}
\end{equation}
The constants and both omitted source tails are explicitly
specified in Section~\ref{sec:canonical}.
\end{theorem}
\begin{proof}
Normalize $F$ and $\widehat F$ as in \eqref{eq:canonical-constants}.
Equation~\eqref{eq:EX} gives $E=E_X$, and their zero populations are
at most $n_X$. Apply Theorem~\ref{thm:pair-transfer} with
\eqref{eq:D2}--\eqref{eq:Ddiag}. Both original cosine fields are
conjugation-symmetric; $D$ is zero-free on the disk. Thus the
folded functionals equal the stated full-plane statistics.
The approximant's zeros in $b<|y|<2b$ retain their actual weights.
For $b=h$, $4n_X(536+164)=2800n_X$ and
$1218+1185/2=3621/2$.
\end{proof}

At the scale \eqref{eq:matched-scale}, the prefactor multiplying
$(\delta/M_X)^{\alpha_X}$ is $O((\log X)^5)$.
The direct log-product allowance \eqref{eq:pair-fourth} instead
has prefactor $O((\log X)^7)$. These are upper-bound costs, not
actual error asymptotics or runtime estimates; both retain the
deteriorating exponent and the enlarged-disk source cost.

Write $K_{\rm pair}=(4n_XD_2+D_{\rm diag})/\pi$.
For any $e>0$, use \eqref{eq:precision-choice} with
$K_\varphi$ replaced by $K_{\rm pair}$ and then
\eqref{eq:tail-allocation}; this gives $E_{\rm pair}\le e/2$.
If a later calculation obtains $\widehat q$ with
$|\widehat q-S(\widehat H_{J,U})|\le E_{\rm eval}$, the precise
condition excluding a plateau pair with $d_0\le d\le b$ is
\begin{equation}\label{eq:pair-consumer}
 \widehat q+E_{\rm eval}+E_{\rm pair}<2(d_0/b)^4 .
\end{equation}
At present the favorable value on the left has not been acquired.

\subsection{Two scale checks}

A center average over an interval of length $L$ containing the
whole pair-window translate dilutes its mutual contribution to
\[
 2(d/b)^4\frac{h}{L}\int_\R w(t)^2\dd t
 =2(d/b)^4\frac{h}{L}\,\frac{643}{462}.
\]
The constant is elementary:
$\int_\R w^2=1+\int_0^1p(t)^2\dd t=1+181/462$.
Thus the fixed-center $1/8$ band margin cannot be carried
unchanged into an averaged theorem.

The scale dependence of the transfer also cannot be discarded.
For $0<d\le1$ consider the local polynomial controls
\begin{equation}\label{eq:pair-control}
 f_\pm(z)=\frac{(z-X)^2\pm d^2}{74}.
\end{equation}
On $\overline D(c,R)$, $|z-X|\le17/2$, so
$\norm{f_\pm}_\infty<1$,
$|f_\pm(c)|\ge15/74$ and
$\norm{f_+-f_-}_\infty=d^2/37$.
Yet $f_+$ has the conjugate pair $X\pm\ii d$ and $f_-$ has
only real roots. Taking $h=b=2d\le1$ gives
$S(f_+)=1/8$ and $S(f_-)=0$ while the field difference tends
to zero. A modulus of continuity independent of the shrinking
test scale is therefore impossible for this generic class.
The controls are local polynomials, not canonical theta fields;
they neither refute Theorem~\ref{thm:canonical-pair} nor decide
a source-specific favorable estimate.

\section{What logarithmic comparison must still control}\label{sec:comparisons}

The approximation theorems are compatible with either a real or a
nonreal local zero configuration. To obtain a favorable-value
estimate one needs information about the canonical field beyond
the error of its approximation. We now make that distinction
quantitative for several natural comparison recipes.

\subsection{The exact signed-Jensen decomposition}

Let $\varphi\ge0$ be a nonzero compact $C^{1,1}$ test in a disk.
Its Laplacian has integral zero. Moreover its positive and
negative parts both have positive mass: otherwise $\Delta\varphi$
would vanish, and a compact harmonic function would be zero.
Set
\begin{equation}\label{eq:signed-measures}
 \beta_\varphi=\int(\Delta\varphi)_+\dd A
             =\int(\Delta\varphi)_-\dd A>0,\qquad
 \dd\mu_\pm=\frac{(\Delta\varphi)_\pm}{\beta_\varphi}\dd A.
\end{equation}
For a nonzero holomorphic $f$ define
\begin{equation}\label{eq:jensen-gaps}
 A_\pm=\int|f|^2\dd\mu_\pm,\qquad
 J_\pm=\log A_\pm-\int\log|f|^2\dd\mu_\pm\ge0.
\end{equation}
All integrals are finite; $A_\pm>0$ because holomorphic zeros have
area measure zero. Jensen's inequality follows by regularization
and a limit if necessary.

\begin{proposition}[Signed-Jensen identity]\label{prop:signed-jensen}
The zero statistic satisfies
\begin{equation}\label{eq:signed-jensen}
 \nu_f(\varphi)=\frac{\beta_\varphi}{4\pi}
       \left[\log\frac{A_+}{A_-}+J_--J_+\right].
\end{equation}
In addition
\begin{equation}\label{eq:quadratic-positive}
 A_+-A_-=\frac4{\beta_\varphi}
                       \int\varphi|f'|^2\dd A\ge0,
\end{equation}
with strict inequality if $f$ is nonconstant.
\end{proposition}
\begin{proof}
Subtract the two equalities in \eqref{eq:jensen-gaps} and
apply \eqref{eq:weak-onepoint}. Integration by parts and
$\Delta|f|^2=4|f'|^2$ give \eqref{eq:quadratic-positive}.
If $f$ is nonconstant, $f'$ cannot vanish on the open set on
which $\varphi>0$.
\end{proof}

The positive sign of the quadratic difference does not control
the signed difference of the Jensen gaps. For a zero-free
nonconstant function, such as $f(z)=\ee^{tz}$ with real $t\ne0$,
the left side of \eqref{eq:signed-jensen} is zero. Thus
$J_--J_+=-\log(A_+/A_-)<0$: even a zero-free field requires
exact cancellation of its quadratic variation.

Nor is $J_--J_+$ always nonpositive. On the unit disk take
$\varphi(z)=(1-|z|^2)_+^2$ and $f(z)=z^2+1/4$.
Writing $u=|z|^2$, one has
$\Delta\varphi=8(2u-1)$, $\beta_\varphi=2\pi$, and the
angular mean of $|f|^2$ is $u^2+1/16$.
Integration over $u>1/2$ and $u<1/2$ respectively gives
\[
 A_+=\frac{17}{24}+\frac1{16}=\frac{37}{48},\qquad
 A_-=\frac1{24}+\frac1{16}=\frac5{48}.
\]
The zeros $\pm\ii/2$ give
$\nu_f(\varphi)=2(1-1/4)^2=9/8$.
Since $9/8>\frac12\log(37/5)$, the gap difference is positive
here. The inequality follows already from the cubic Taylor lower
bound for $\ee^{9/4}>37/5$. This is a control on the generic
argument, not a canonical theta counterexample.

The quadratic quantities themselves have a literal source formula.
For $f=F/M_X$, substitute \eqref{eq:source} and
\eqref{eq:canonical-F} into $A_\pm$:
\begin{align}
 A_\pm&=M_X^{-2}\iint_{(0,\infty)^2}
               \Phi(u)\Phi(v)\mathcal K_\pm(u,v)\dd u\dd v,\label{eq:quadratic-source}\\
 \mathcal K_\pm(u,v)&=
       \int |D(z)|^2\cos(zu)\overline{\cos(zv)}\dd\mu_\pm(z).\notag
\end{align}
Fubini follows from the complete source majorant on the compact
disk. Each kernel defines a nonnegative quadratic form, but this
does not compare the logarithmic gaps.

With $g=\widehat F/M_X$ and $\varepsilon=\delta/M_X$,
$|f|,|g|\le1$ imply $|A_\pm-\widehat A_\pm|\le2\varepsilon$.
If $\widehat A_\pm>2\varepsilon$, the corresponding log error is at
most $2\varepsilon/(\widehat A_\pm-2\varepsilon)$.
Also \eqref{eq:EX} gives
\[
 \left|\int\log|f|^2\dd\mu_\pm-
               \int\log|g|^2\dd\mu_\pm\right|
 \le\frac{2\norm{\Delta\varphi}_\infty}{\beta_\varphi}E_X.
\]
These estimates acquire the separate quantities. They do not
bound their combination favorably.

\subsection{A sharp local distance from harmonic functions}

If a comparison field is zero-free, its logarithm is harmonic.
The following estimate measures the unavoidable absolute
logarithmic cost of using such a baseline near a zero.

\begin{theorem}[Sharp harmonic-distance bound]\label{thm:harmonic-distance}
Let $F$ be nonzero and holomorphic on $D(a,\ell)$, with a zero
of multiplicity $m\ge1$ at $a$. For every real harmonic function
$v$ on that disk,
\begin{equation}\label{eq:harmonic-distance}
 \int_{D(a,\ell)}|\log|F|-v|\dd A
 \ge\frac{\pi\log2}{2}\,m\ell^2 .
\end{equation}
The constant is sharp.
\end{theorem}
\begin{proof}
For $0<\varrho<\ell$, let $t=|z-a|$ and define
\[
 \Psi_\varrho(t)=
 \begin{cases}
 (\varrho^2\log2-t^2)/4,
                      &0\le t\le\varrho/\sqrt2,\\[2pt]
 (t^2-\varrho^2)/4-(\varrho^2/2)\log(t/\varrho),
                      &\varrho/\sqrt2\le t\le\varrho,\\[2pt]
 0,                   &t\ge\varrho.
 \end{cases}
\]
The value and first derivative agree at both junctions.
This is a nonnegative compact $C^{1,1}$ function, decreasing to
zero at $\varrho$, with
\[
 \Delta\Psi_\varrho=
 \begin{cases}-1,&t<\varrho/\sqrt2,\\
                1,&\varrho/\sqrt2<t<\varrho,\\
                0,&t>\varrho,\end{cases}
 \qquad
 \Psi_\varrho(0)=\frac{\varrho^2\log2}{4}.
\]
In particular $\norm{\Delta\Psi_\varrho}_\infty=1$, and there
are no boundary delta terms. The weak zero formula and harmonicity
give
\[
 \int(\log|F|-v)\Delta\Psi_\varrho\dd A
 =2\pi\sum_{F(z)=0}m_z\Psi_\varrho(z)
 \ge2\pi m\Psi_\varrho(0).
\]
Absolute values bound the left side by the $L^1$ distance.
Letting $\varrho\uparrow\ell$ proves
\eqref{eq:harmonic-distance}, also when that distance is infinite.

For sharpness take $F(z)=(z-a)^m$ and the constant harmonic
function $v=m\log(\ell/\sqrt2)$. With $u=|z-a|^2/\ell^2$,
the distance equals
\[
 \frac{\pi m\ell^2}{2}\int_0^1|\log(2u)|\dd u
 =\frac{\pi m\ell^2\log2}{2}.
\]
This matches the lower bound.
\end{proof}

Consider the precise recipe in which one estimates
$S_\varphi(F)=\nu_F(\varphi)$ by subtracting a harmonic $v$ and
paying the full-disk absolute allowance
\[
 B\ge\norm{\log|F|-v}_{L^1(D(c,r))},\qquad
 |S_\varphi(F)|\le K_\varphi B.
\]
If $X$ itself is a real zero of multiplicity $m$, the disk
$D(X,1/6)$ lies in $D(c,r)$, because $r-4=1/6$.
Theorem~\ref{thm:harmonic-distance} forces
\begin{equation}\label{eq:baseline-floor}
 B\ge\frac{\pi m\log2}{72}.
\end{equation}
For the scalar tests on their plateau, just above the real axis,
\[
 \Delta\varphi_Q=4,\qquad
 \Delta\varphi_\eta=
        \frac{4\eta(1-\eta y^2)}{(1+\eta y^2)^2}
                  \longrightarrow4\eta .
\]
Thus this recipe's allowances satisfy
\begin{equation}\label{eq:allowance-floors}
 K_QB\ge\frac{m\log2}{36},\qquad
 K_\eta B\ge\frac{\eta m\log2}{36}.
\end{equation}
They cannot pass the strict scalar consumers for
$d^2\le m\log2/72$, since
$2\log(1+\eta d^2)\le2\eta d^2$.
These are floors on the \emph{allowance}, not lower bounds on
the true transverse mass, which may be zero.
The excluded class is the harmonic-baseline comparison with this
full-disk absolute payment in a real-zero-centered window.
No conclusion follows about all weighted, signed, or real-divisor
comparisons.

\subsection{Keeping real-axis vanishing by regularization}

Quadratic vanishing of $\varphi$ does not make
$|\Delta\varphi|$ vanish at the axis. In particular the immediate
log-error bound is
$(2\pi)^{-1}\int|\Delta\varphi|\,|e|\dd A$,
not a $y^2$-weighted log norm. On a plateau,
$\Delta\varphi_Q=4$ for small positive $y$.
If arbitrary nonnegative residuals are concentrated in
$0<y<\epsilon$, a bound by
$C\int y^\beta |e|\dd A$, $\beta>0$, would require
$C\gtrsim\epsilon^{-\beta}$.
This only rules out the unqualified norm substitution; those
residuals are not asserted to be holomorphic log differences.

A legitimate way to retain the original test is to differentiate
a regularized logarithm. For $0<\tau\le1$ set
\[
 u_\tau(f)=\frac12\log(|f|^2+\tau^2),\qquad
 T_{\varphi,\tau}(f)=\frac1\pi
       \int\varphi\,\frac{\tau^2|f'|^2}
                            {(|f|^2+\tau^2)^2}\dd A .
\]
Holomorphy gives
$\Delta u_\tau=2\tau^2|f'|^2/(|f|^2+\tau^2)^2$, hence
$T_{\varphi,\tau}=(2\pi)^{-1}\int u_\tau\Delta\varphi\dd A$.
For $\varphi\ge0$ this integral is nonnegative. As $\tau\downarrow0$,
$u_\tau\downarrow\log|f|$ in local $L^1$, by domination with
$u_1-\log|f|$. Therefore
\begin{equation}\label{eq:regularized-limit}
 T_{\varphi,\tau}(f)\longrightarrow\nu_f(\varphi).
\end{equation}
Real multiple zeros receive their zero test weight in this limit;
nonreal zeros retain their positive weight.

The sublevel estimate \eqref{eq:sublevel} supplies a quantitative
bias when $n\ge2$. Write $L=L(A,n)$. Layer cake gives
\begin{align}
 \int(u_\tau(f)-\log|f|)\dd A
 &\le L\tau^{2/n}
       \int_0^\infty(\ee^{2v}-1)^{-1/n}\dd v\notag\\
 &\le 2nL\tau^{2/n}.\label{eq:regularization-bias}
\end{align}
Indeed $u_\tau-\log|f|>v$ means
$|f|<\tau(\ee^{2v}-1)^{-1/2}$.
On $0<v<1$, use $\ee^{2v}-1\ge2v$ to bound the last integral
by two; on $v\ge1$, use $\ee^{2v}-1\ge\ee^v$ to bound the
tail by $n$. Then $n+2\le2n$.
The canonical $n_X$ is at least two. Consequently
\[
 |T_{\varphi,\tau}(g)-\nu_g(\varphi)|
 \le B_{\varphi,\tau}
 :=2K_\varphi n_XL(A_X,n_X)\tau^{2/n_X}.
\]
This is a smoothing bias for one logarithm; its exponent
$2/n_X$ does not replace the source-transfer exponent
$2/(n_X+2)$.

The complete possible consumer would have to bound
\begin{equation}\label{eq:regularized-consumer}
 T_{\varphi,\tau}(g)+B_{\varphi,\tau}
       +K_\varphi E_X+E_{\rm eval}
       <\hbox{the declared defect margin}.
\end{equation}
The first term requires joint information about small values and
derivatives. The generic Cauchy bound from $R-r=1/3$ is
$|g'|\le3$, and yields only
$T_{\varphi,\tau}(g)\le9(\pi\tau^2)^{-1}\int\varphi\dd A$.
This allowance grows as the smoothing bias shrinks. The present
complete-source upper bounds therefore do not close this recipe.
This conclusion is scoped to those available inputs, not to
every signed or weighted logarithmic method.

\subsection{Removing a real divisor and accounting for what remains}

Let $P_R(z)=\prod_\ell(z-a_\ell)^{m_\ell}$ with real $a_\ell$.
For either scalar test, $\nu_{P_R}(\varphi)=0$.
For any harmonic $v$ one consequently has the identity
\begin{equation}\label{eq:real-divisor}
 \nu_F(\varphi)=\frac1{2\pi}\int\Delta\varphi
     \bigl(\log|F|-\log|P_R|-v\bigr)\dd A.
\end{equation}
It remains valid when $P_R$ is not a divisor of $F$: the
corresponding real poles have zero test weight.
To conclude that $F/P_R$ is holomorphic, however, every selected
root and its multiplicity must be justified independently.
Even removal of all real roots in a compact domain leaves any
nonreal roots in the quotient.

This representation does improve on the harmonic class for a
purely real divisor. If $F=(z-a)^m\ee^q$, with real $a$ and
holomorphic $q$, take $P_R=(z-a)^m$ and $v=\Re q$; the residual
is exactly zero. The harmonic allowance floor therefore does
not extend to this class.
Two costs remain. A field error cannot be divided by $P_R$
at its roots without shared vanishing or controlled deflation:
$F=z-a$, $\widehat F=z-a+\epsilon$ is an immediate example.
Also, the finite theta source has no acquired real-rootedness
theorem, so it cannot simply be declared an all-real comparator.

A concrete independent alternative would be a matched integer
count. Let a Jordan domain $\Omega$ contain the whole test
support, have closure within $D(c,r)$, and satisfy
$F\ne0$ on its boundary. Suppose rigorous integers give
\[
 N_\Omega=\frac1{2\pi\ii}\int_{\partial\Omega}\frac{F'}F\dd z
          \le U,\qquad
 \hbox{certified real multiplicity in }\Omega\ge M_R.
\]
The argument principle makes $N_\Omega$ the total multiplicity.
Thus $R_\Omega=U-M_R\ge0$ bounds all unaccounted zeros,
including every relevant nonreal zero. The canonical strip
then gives the valid coarse estimates
\begin{equation}\label{eq:count-only}
 Q\le2R_\Omega,\qquad
 \log Z(\eta)\le2\log(1+\eta)R_\Omega .
\end{equation}
If $R_\Omega\ge1$, these particular allowances are at least
the one-pair margin for every $0<d\le1$. The count-only recipe
therefore needs $U=M_R$ to pass the strict scalar consumer.
In that case all zeros in the domain are accounted for as real,
and the statistic is zero directly. Matching counts, a positive
boundary modulus, and real-axis sign or multiplicity certificates
have not been supplied here. Cauchy's estimate gives derivative
approximation error at most $3\delta$ on $D(c,r)$; it gives
no lower bound for $|F|$ on a proposed contour.

Finally, a common real divisor and a safe anchor cannot by
themselves force reality at fixed positive approximation tolerance.
The local controls
\begin{equation}\label{eq:real-divisor-control}
 f_\pm(z)=\frac{(z-X)^2\pm d^2}{5(-16\pm d^2)},
 \qquad 0<d\le1,
\end{equation}
have $f_\pm(c)=1/5$ and
$\norm{f_\pm}_\infty\le293/300<1$ on $D(c,R)$.
Direct subtraction yields
\[
 f_+-f_-=
 -\frac{2d^2((z-X)^2+16)}{5(256-d^4)},\qquad
 \norm{f_+-f_-}_\infty\le\frac{353}{2550}d^2.
\]
Their transverse masses are $2d^2$ and zero. Multiplying both
by a fixed real polynomial divided by its disk supremum
preserves a common known real divisor, an upper bound one,
and a common positive anchor modulus; all added real zeros
remain invisible to the test. These controls lack the canonical
theta source identity. They show why fixed generic local
information is insufficient, without excluding a theta-specific
estimate, scale-dependent precision, or a valid finite count
certificate.

\section{Evidence status and supporting calculations}\label{app:status}

This working manuscript was drafted for Edward Baker on
8 November 2026 with substantial OpenAI Codex LLM assistance.
The current recorded source does not expose the actual backend
model or reasoning effort; neither is inferred. No model switch,
separate reviewer, or independent specialist assessment occurred
in the historical manuscript-rewrite turn. A subsequent separate-session
E07 assessment used three read-only inherited-model subchecks and root
synthesis, with one bounded correction. This is internal LLM assessment,
not independent specialist or verified cross-model validation.

The source correspondence is as follows. The cited internal
notes contain the dated original arguments and their assessment
records; the manuscript reorders and derives their relevant
results in a single notation.

\begin{center}\small
\begin{tabularx}{\textwidth}{@{}>{\raggedright\arraybackslash}p{0.27\textwidth}Y@{}}
\toprule
Manuscript content & Source and present evidence status\\
\midrule
Full Gaussian factorial hierarchy, weighted observation law,
center projections and covariance &
Note09 and this manuscript \cite{OCGaussian}; accepted in the
subsequent internal E07 assessment, with roman subscripts repaired. Standard framework distinguished from canonical conclusions.\\
Two-sided window and uniform full-pair source approximation &
Note09 and Theorem~\ref{thm:full-pair}; derived from accepted E02
inputs; accepted internally in E07, with an explicit localization-width
floor for common-cutoff uniform density tolerance. No canonical profile estimate.\\
Finite reflection and configuration algebra &
Note01, incorporated in the E02 assessment
\cite{OCFinite,OCReviewE02}. Accepted internally after the explicit
window-support correction for the count corollary.\\
Local logarithm lemma and complete one-point approximation &
Notes02--03 and E02 \cite{OCLog,OCConsumer,OCReviewE02}.
Accepted in that internal LLM assessment; no favorable canonical
value was established.\\
Signed-Jensen and harmonic-distance results &
B02 Note04 \cite{OCHarmonic}. Proved in-note; proof-author
self-check only. E03 remains pending.\\
Regularization, real-divisor and matched-count discussion &
B03 comparison note \cite{OCRealDivisor}. Proved in-note;
author checking only. Its two completion attempts are parked
on the currently acquired inputs.\\
Marked-pair transfer and explicit constants &
B04 Note06 and author check04 \cite{OCPair,OCB04Check}.
Proved in-note; proof-author self-check only. E04 remains pending.\\
BGST reconciliation and arithmetic projection &
B04 Note05 and the separate bounded source audit
\cite{OCArithmetic,OCBGSTAudit}; primary version
\cite{BGST2026PairI}. Selected source/import checking does not
assess the B04 pair-transfer proof.\\
\bottomrule
\end{tabularx}
\end{center}

The new derivations and the elementary source-normalization calculation, expanded
intermediate derivations and source crosswalk in this assembly
had a same-assistant integration check at E06. E07 subsequently assessed
only the full-function claims and their E02 interface as stated above;
the earlier B02/B04 supporting claims were not assessed. The theorem labels
identify written mathematical claims; they do not erase the
different assessment levels above. In particular, no source is
promoted by being placed in a manuscript.

The investigation contains small exact-rational bookkeeping
records for constants, cutoff integrals, and multiplicity
coefficients. They support transcription checks but replace
none of the analytical arguments. The corrected optional
averaging constant is $643/462$, as derived in
Appendix~\ref{sec:pairs}; an earlier working arithmetic slip was
recorded and corrected in the B04 author check.
No canonical zero sampling, source quadrature campaign,
floating correlation fit, or numerical evidence for RH is used
in this manuscript. The build record, source hashes, frozen-input
checks, and visual inspection establish artifact provenance and
production quality only.

The rewrite makes full marginals primary and moves the earlier scalar
analysis into supporting appendices. It supersedes the previous
manuscript's emphasis, not its dated production record or the frozen
assessment inputs. The new author check records exact support-buffer
and weighted-tuple checks; none is a substitute for proof assessment.

\bibliographystyle{plainurl}
\bibliography{localized_zero_correlations_sources}
\end{document}
